% Generated human-review packet template.  Source records, Lean statements,
% and raw-source-to-Spec screening fields are substituted by
% scripts/review_dashboard_packet.py.
% Do not hand-edit generated packets; annotate the PDF or reconcile notes into
% the paper-local human-review log.
\documentclass[10pt]{article}
\usepackage[margin=0.43in,footskip=0.2in]{geometry}
\usepackage{fontspec}
% Use a Unicode-complete main face as well as a Unicode-complete code face.
% Reviewer reasons and source labels can contain exact Greek symbols outside
% verbatim blocks; silently dropping one would corrupt the review packet.
\setmainfont{DejaVu Serif}
% The broad DejaVu Sans Unicode coverage preserves Lean's mathematical glyphs
% (including U+2216 set difference and superscript identifiers) in verbatim
% review blocks.  Its proportional metrics are preferable to silently missing
% exact semantic text from a narrower monospace font.
\setmonofont{DejaVu Sans}
\usepackage{hyperref}
\usepackage{amssymb}
\usepackage{fvextra}
\usepackage{enumitem}
\setlength{\parindent}{0pt}
\setlength{\parskip}{0.15em}
\linespread{0.92}
\hypersetup{colorlinks=true,linkcolor=blue,urlcolor=blue,breaklinks=true}
\DefineVerbatimEnvironment{ReviewVerbatim}{Verbatim}{breaklines=true,breakanywhere=true,fontsize=\scriptsize}
\newcommand{\reviewerbox}{%
  \par\noindent\fbox{\parbox{0.96\linewidth}{\rule{0pt}{0.38in}}}\par
}
\newcommand{\reviewmatch}[1]{%
  \par\noindent\textbf{Reviewer check:} $\square$\; #1\par
  \noindent\emph{If left unchecked, explain the discrepancy or uncertainty below.}\par
}

\begin{document}
\fontsize{9}{10}\selectfont

\begin{center}
{\LARGE Human review packet: DSWG24DiscretizationBias}\\[0.35em]
\small Generated from byte-pinned source excerpts, the expanded PaperInterface
specifications, and hash-bound raw-source-to-Spec screening records.
\end{center}

\noindent\textbf{Paper:} DSWG24DiscretizationBias\\
\textbf{Paper id:} \texttt{DSWG24DiscretizationBias}\\
\textbf{Source version:} \parbox[t]{0.76\linewidth}{\raggedright arXiv:2405.16762; byte-pinned publication source record}\\
\textbf{Source record:} \texttt{audit/paper\_statement\_map.json}\\
\textbf{Rows:} 9\\
\textbf{Generated:} 2026-09-07\\
\textbf{Source URL:} \href{https://arxiv.org/pdf/2405.16762}{official source}





\section*{How to use this packet}
For each source claim, compare the exact source excerpts with the expanded
PaperInterface specification. Mark up this PDF or fill the blank reviewer box in the TeX source;
return the annotated file for reconciliation into the paper-local human-review
log.

\section*{Open the interactive dashboard}
The interactive dashboard is an optional alternative to using this PDF; it
presents the same review surface rather than an additional review requirement.
After forking and cloning (or pulling) AppliedModelingLib, run the following from the
repository root. The cache command is needed only the first time in a new clone.
\begin{ReviewVerbatim}
cd <your-repository-checkout>
lake exe cache get
python3 scripts/review_dashboard.py --paper DSWG24DiscretizationBias --serve
\end{ReviewVerbatim}
Open \url{http://127.0.0.1:8765/} in a browser and keep that terminal running
while reviewing. The dashboard records saved annotations in this paper's local
reviewer-owned review record.

\section*{Contents}
\begin{itemize}[leftmargin=1.2em]
\item \hyperlink{governing-declaration-section-5ef5ef0364b6939c}{Governing Lean declarations (49; supporting context)}
\item \hyperlink{paper-prerequisite-section-5ef5ef0364b6939c}{Paper-specific semantic prerequisites (11; not paper claims)}
\begin{itemize}[leftmargin=1.2em]
\item \hyperlink{paper-prerequisite-c793e07e52a48893}{DSWG24\allowbreak{}Discretization\allowbreak{}Bias.\allowbreak{}Proof\allowbreak{}Bridge.\allowbreak{}bayes\allowbreak{}Optimal}
\item \hyperlink{paper-prerequisite-c94cbc772446ffec}{DSWG24\allowbreak{}Discretization\allowbreak{}Bias.\allowbreak{}Proof\allowbreak{}Bridge.\allowbreak{}calibrated}
\item \hyperlink{paper-prerequisite-46e579d6fcd53814}{DSWG24\allowbreak{}Discretization\allowbreak{}Bias.\allowbreak{}Proof\allowbreak{}Bridge.\allowbreak{}is\allowbreak{}Aggregate\allowbreak{}First\allowbreak{}Argmax}
\item \hyperlink{paper-prerequisite-2e16c9b0a6f0db1c}{DSWG24\allowbreak{}Discretization\allowbreak{}Bias.\allowbreak{}Proof\allowbreak{}Bridge.\allowbreak{}is\allowbreak{}Argmax\allowbreak{}Rule}
\item \hyperlink{paper-prerequisite-536b545abe995da0}{DSWG24\allowbreak{}Discretization\allowbreak{}Bias.\allowbreak{}Proof\allowbreak{}Bridge.\allowbreak{}is\allowbreak{}Independent\allowbreak{}Randomized\allowbreak{}Rule}
\item \hyperlink{paper-prerequisite-45b940eee9bfbd29}{DSWG24\allowbreak{}Discretization\allowbreak{}Bias.\allowbreak{}Proof\allowbreak{}Bridge.\allowbreak{}is\allowbreak{}Independent\allowbreak{}Rule}
\item \hyperlink{paper-prerequisite-29d817bb87f95d82}{DSWG24\allowbreak{}Discretization\allowbreak{}Bias.\allowbreak{}Proof\allowbreak{}Bridge.\allowbreak{}is\allowbreak{}Thompson\allowbreak{}Sampling\allowbreak{}Rule}
\item \hyperlink{paper-prerequisite-f1caeaac865e715b}{DSWG24\allowbreak{}Discretization\allowbreak{}Bias.\allowbreak{}Proof\allowbreak{}Bridge.\allowbreak{}is\allowbreak{}Tie\allowbreak{}Broken\allowbreak{}Argmax\allowbreak{}Rule}
\item \hyperlink{paper-prerequisite-ce0ba7b00b93c0fc}{DSWG24\allowbreak{}Discretization\allowbreak{}Bias.\allowbreak{}Proof\allowbreak{}Bridge.\allowbreak{}maximizes\allowbreak{}Equation1}
\item \hyperlink{paper-prerequisite-b11323803f325211}{DSWG24\allowbreak{}Discretization\allowbreak{}Bias.\allowbreak{}Proof\allowbreak{}Bridge.\allowbreak{}posterior\allowbreak{}Simplex}
\item \hyperlink{paper-prerequisite-e9ff5dcccf634407}{DSWG24\allowbreak{}Discretization\allowbreak{}Bias.\allowbreak{}Proof\allowbreak{}Bridge.\allowbreak{}reference\allowbreak{}Simplex\allowbreak{}At}
\end{itemize}
\item \hyperlink{source-section-f10b5f68e4acf3dc}{Source claims (9)}
\begin{itemize}[leftmargin=1.2em]
\item \hyperlink{source-claim-1966e78967344232}{theorem1ii\_\allowbreak{}perfect\_\allowbreak{}classifier\_\allowbreak{}zero\_\allowbreak{}bias\allowbreak{}Spec}
\item \hyperlink{source-claim-61c8bcc1f7628e92}{theorem2iii\_\allowbreak{}randomized\_\allowbreak{}non\_\allowbreak{}argmax\_\allowbreak{}not\_\allowbreak{}pareto\allowbreak{}Spec}
\item \hyperlink{source-claim-1356bc033e8e977b}{theorem2iii\_\allowbreak{}randomized\_\allowbreak{}weighted\_\allowbreak{}objective\_\allowbreak{}maximizer\_\allowbreak{}agrees\_\allowbreak{}argmax\allowbreak{}Spec}
\item \hyperlink{source-claim-4662dc8e8bb960d7}{nontrivial\_\allowbreak{}reference\_\allowbreak{}family\_\allowbreak{}definition\allowbreak{}Spec}
\item \hyperlink{source-claim-861b407dcf52f582}{theorem1i\_\allowbreak{}no\_\allowbreak{}information\_\allowbreak{}bias\allowbreak{}Spec}
\item \hyperlink{source-claim-d8ad786b7cc58022}{theorem1iii\_\allowbreak{}argmax\_\allowbreak{}bias\_\allowbreak{}le\_\allowbreak{}mae\allowbreak{}Spec}
\item \hyperlink{source-claim-d6b9878258465151}{theorem1iii\_\allowbreak{}tight\_\allowbreak{}binary\_\allowbreak{}example\allowbreak{}Spec}
\item \hyperlink{source-claim-82ddd45159ec7ca3}{theorem2i\_\allowbreak{}joint\_\allowbreak{}rule\_\allowbreak{}exists\allowbreak{}Spec}
\item \hyperlink{source-claim-cc8ba41c41752422}{theorem2ii\_\allowbreak{}argmax\_\allowbreak{}accuracy\_\allowbreak{}maximizing\allowbreak{}Spec}
\end{itemize}
\end{itemize}

\clearpage
\hypertarget{governing-declaration-section-5ef5ef0364b6939c}{}
\section*{Governing Lean declarations}
These exact graph-bound declaration bodies are retained by the displayed specifications or reviewed prerequisites. They are supporting context and do not add source-result rows or semantic judgments.
\hypertarget{governing-declaration-06698b1aac1f5804}{}
\subsection*{\small\ttfamily\raggedright Applied\allowbreak{}Modeling\allowbreak{}Lib.\allowbreak{}Decision.\allowbreak{}Is\allowbreak{}Pointwise\allowbreak{}Max}
\textbf{Declaration location:} reusable library
\paragraph{Lean declaration body}
\begin{ReviewVerbatim}
type:
{ι : Type u_1} → {α : Type u_2} → (ι → α → ℝ) → (ι → α) → Prop

value:
fun {ι : Type u_1} {α : Type u_2} (score : ι → α → ℝ) (choose : ι → α) =>
  ∀ (i : ι) (a : α), score i a ≤ score i (choose i)
\end{ReviewVerbatim}


\hypertarget{governing-declaration-eb5e05858b12cdd2}{}
\subsection*{\small\ttfamily\raggedright Applied\allowbreak{}Modeling\allowbreak{}Lib.\allowbreak{}Decision.\allowbreak{}Monotone\allowbreak{}Expectation}
\textbf{Declaration location:} reusable library
\paragraph{Lean declaration body}
\begin{ReviewVerbatim}
type:
{ω : Type u_1} → ((ω → ℝ) → ℝ) → Prop

value:
fun {ω : Type u_1} (expect : (ω → ℝ) → ℝ) => ∀ (f g : ω → ℝ), (∀ (x : ω), f x ≤ g x) → expect f ≤ expect g
\end{ReviewVerbatim}


\hypertarget{governing-declaration-443446486047239c}{}
\subsection*{\small\ttfamily\raggedright Applied\allowbreak{}Modeling\allowbreak{}Lib.\allowbreak{}Decision.\allowbreak{}Finite\allowbreak{}Linear\allowbreak{}Expectation}
\textbf{Declaration location:} reusable library
\paragraph{Lean declaration body}
\begin{ReviewVerbatim}
type:
{ω : Type u_1} → ((ω → ℝ) → ℝ) → Prop

value:
fun {ω : Type u_1} (expect : (ω → ℝ) → ℝ) =>
  AppliedModelingLib.Decision.MonotoneExpectation expect ∧
    (∀ (c : ℝ) (f : ω → ℝ), (expect fun (x : ω) => c * f x) = c * expect f) ∧
      ∀ (f g : ω → ℝ), (expect fun (x : ω) => f x + g x) = expect f + expect g
\end{ReviewVerbatim}

\paragraph{Direct retained dependencies}
\begin{itemize}\raggedright
\item \hyperlink{governing-declaration-eb5e05858b12cdd2}{Applied\allowbreak{}Modeling\allowbreak{}Lib.\allowbreak{}Decision.\allowbreak{}Monotone\allowbreak{}Expectation}
\end{itemize}
\hypertarget{governing-declaration-3088d80cd010ba5f}{}
\subsection*{\small\ttfamily\raggedright Applied\allowbreak{}Modeling\allowbreak{}Lib.\allowbreak{}Decision.\allowbreak{}average\allowbreak{}Score}
\textbf{Declaration location:} reusable library
\paragraph{Lean declaration body}
\begin{ReviewVerbatim}
type:
{ι : Type u_1} → {α : Type u_2} → [Fintype ι] → (ι → α → ℝ) → (ι → α) → ℝ

value:
fun {ι : Type u_1} {α : Type u_2} [Fintype ι] (score : ι → α → ℝ) (choose : ι → α) =>
  (∑ i : ι, score i (choose i)) / ↑(Fintype.card ι)
\end{ReviewVerbatim}


\hypertarget{governing-declaration-0cd266624e93f607}{}
\subsection*{\small\ttfamily\raggedright Applied\allowbreak{}Modeling\allowbreak{}Lib.\allowbreak{}Decision.\allowbreak{}dataset\allowbreak{}Accuracy}
\textbf{Declaration location:} reusable library
\paragraph{Lean declaration body}
\begin{ReviewVerbatim}
type:
{ι : Type u_1} → {α : Type u_2} → [Fintype ι] → [DecidableEq α] → (ι → α) → (ι → α) → ℝ

value:
fun {ι : Type u_1} {α : Type u_2} [Fintype ι] [DecidableEq α] (truth decision : ι → α) =>
  AppliedModelingLib.Decision.averageScore (fun (i : ι) (a : α) => if a = truth i then 1 else 0) decision
\end{ReviewVerbatim}

\paragraph{Direct retained dependencies}
\begin{itemize}\raggedright
\item \hyperlink{governing-declaration-3088d80cd010ba5f}{Applied\allowbreak{}Modeling\allowbreak{}Lib.\allowbreak{}Decision.\allowbreak{}average\allowbreak{}Score}
\end{itemize}
\hypertarget{governing-declaration-61110afe97ed1775}{}
\subsection*{\small\ttfamily\raggedright Applied\allowbreak{}Modeling\allowbreak{}Lib.\allowbreak{}Decision.\allowbreak{}expected\allowbreak{}Decision\allowbreak{}Accuracy}
\textbf{Declaration location:} reusable library
\paragraph{Lean declaration body}
\begin{ReviewVerbatim}
type:
{ω : Type u_1} →
  {σ : Type u_2} →
    {ι : Type u_3} →
      {α : Type u_4} → [Fintype ι] → [DecidableEq α] → ((ω → ℝ) → ℝ) → (ω → σ) → (ω → ι → α) → (σ → ι → α) → ℝ

value:
fun {ω : Type u_1} {σ : Type u_2} {ι : Type u_3} {α : Type u_4} [Fintype ι] [DecidableEq α] (expect : (ω → ℝ) → ℝ)
    (obs : ω → σ) (truth : ω → ι → α) (rule : σ → ι → α) =>
  expect fun (x : ω) => AppliedModelingLib.Decision.datasetAccuracy (truth x) (rule (obs x))
\end{ReviewVerbatim}

\paragraph{Direct retained dependencies}
\begin{itemize}\raggedright
\item \hyperlink{governing-declaration-0cd266624e93f607}{Applied\allowbreak{}Modeling\allowbreak{}Lib.\allowbreak{}Decision.\allowbreak{}dataset\allowbreak{}Accuracy}
\end{itemize}
\hypertarget{governing-declaration-71856cc0344b281a}{}
\subsection*{\small\ttfamily\raggedright Applied\allowbreak{}Modeling\allowbreak{}Lib.\allowbreak{}Decision.\allowbreak{}expected\allowbreak{}Objective}
\textbf{Declaration location:} reusable library
\paragraph{Lean declaration body}
\begin{ReviewVerbatim}
type:
{ω : Type u_1} → {σ : Type u_2} → {β : Type u_3} → ((ω → ℝ) → ℝ) → (ω → σ) → (σ → β → ℝ) → (σ → β) → ℝ

value:
fun {ω : Type u_1} {σ : Type u_2} {β : Type u_3} (expect : (ω → ℝ) → ℝ) (obs : ω → σ) (objective : σ → β → ℝ)
    (rule : σ → β) =>
  expect fun (x : ω) => objective (obs x) (rule (obs x))
\end{ReviewVerbatim}


\hypertarget{governing-declaration-28219918dd3306b5}{}
\subsection*{\small\ttfamily\raggedright Applied\allowbreak{}Modeling\allowbreak{}Lib.\allowbreak{}Learning.\allowbreak{}Prediction.\allowbreak{}Model}
\textbf{Declaration location:} reusable library
\paragraph{Lean declaration body}
\begin{ReviewVerbatim}
type:
Type u_1 → Type u_2 → Type (max u_1 u_2)

value:
fun (X : Type u_1) (Y : Type u_2) => X → Y
\end{ReviewVerbatim}


\hypertarget{governing-declaration-48c36d87443b2ef9}{}
\subsection*{\small\ttfamily\raggedright Applied\allowbreak{}Modeling\allowbreak{}Lib.\allowbreak{}Learning.\allowbreak{}Prediction.\allowbreak{}Multiclass\allowbreak{}Score}
\textbf{Declaration location:} reusable library
\paragraph{Lean declaration body}
\begin{ReviewVerbatim}
type:
Type u_1 → Type u_2 → Type (max u_1 u_2)

value:
fun (X : Type u_1) (Y : Type u_2) => AppliedModelingLib.Learning.Prediction.Model X (Y → ℝ)
\end{ReviewVerbatim}

\paragraph{Direct retained dependencies}
\begin{itemize}\raggedright
\item \hyperlink{governing-declaration-28219918dd3306b5}{Applied\allowbreak{}Modeling\allowbreak{}Lib.\allowbreak{}Learning.\allowbreak{}Prediction.\allowbreak{}Model}
\end{itemize}
\hypertarget{governing-declaration-d9f4aa229cbce6df}{}
\subsection*{\small\ttfamily\raggedright Applied\allowbreak{}Modeling\allowbreak{}Lib.\allowbreak{}Learning.\allowbreak{}Prediction.\allowbreak{}Is\allowbreak{}Argmax\allowbreak{}Rule}
\textbf{Declaration location:} reusable library
\paragraph{Lean declaration body}
\begin{ReviewVerbatim}
type:
{X : Type u_1} →
  {Y : Type u_2} →
    AppliedModelingLib.Learning.Prediction.MulticlassScore X Y → AppliedModelingLib.Learning.Prediction.Model X Y → Prop

value:
fun {X : Type u_1} {Y : Type u_2} (q : AppliedModelingLib.Learning.Prediction.MulticlassScore X Y)
    (rule : AppliedModelingLib.Learning.Prediction.Model X Y) =>
  AppliedModelingLib.Decision.IsPointwiseMax q rule
\end{ReviewVerbatim}

\paragraph{Direct retained dependencies}
\begin{itemize}\raggedright
\item \hyperlink{governing-declaration-06698b1aac1f5804}{Applied\allowbreak{}Modeling\allowbreak{}Lib.\allowbreak{}Decision.\allowbreak{}Is\allowbreak{}Pointwise\allowbreak{}Max}
\item \hyperlink{governing-declaration-28219918dd3306b5}{Applied\allowbreak{}Modeling\allowbreak{}Lib.\allowbreak{}Learning.\allowbreak{}Prediction.\allowbreak{}Model}
\item \hyperlink{governing-declaration-48c36d87443b2ef9}{Applied\allowbreak{}Modeling\allowbreak{}Lib.\allowbreak{}Learning.\allowbreak{}Prediction.\allowbreak{}Multiclass\allowbreak{}Score}
\end{itemize}
\hypertarget{governing-declaration-d1592e772e7c7e3d}{}
\subsection*{\small\ttfamily\raggedright Applied\allowbreak{}Modeling\allowbreak{}Lib.\allowbreak{}Learning.\allowbreak{}Prediction.\allowbreak{}Is\allowbreak{}Simplex\allowbreak{}Valued}
\textbf{Declaration location:} reusable library
\paragraph{Lean declaration body}
\begin{ReviewVerbatim}
type:
{X : Type u_1} → {Y : Type u_2} → [Fintype Y] → AppliedModelingLib.Learning.Prediction.MulticlassScore X Y → Prop

value:
fun {X : Type u_1} {Y : Type u_2} [Fintype Y] (q : AppliedModelingLib.Learning.Prediction.MulticlassScore X Y) =>
  (∀ (x : X), ∑ y : Y, q x y = 1) ∧ (∀ (x : X) (y : Y), 0 ≤ q x y) ∧ ∀ (x : X) (y : Y), q x y ≤ 1
\end{ReviewVerbatim}

\paragraph{Direct retained dependencies}
\begin{itemize}\raggedright
\item \hyperlink{governing-declaration-48c36d87443b2ef9}{Applied\allowbreak{}Modeling\allowbreak{}Lib.\allowbreak{}Learning.\allowbreak{}Prediction.\allowbreak{}Multiclass\allowbreak{}Score}
\end{itemize}
\hypertarget{governing-declaration-8e4b4389a4f6454b}{}
\subsection*{\small\ttfamily\raggedright Applied\allowbreak{}Modeling\allowbreak{}Lib.\allowbreak{}pmf\allowbreak{}Exp}
\textbf{Declaration location:} reusable library
\paragraph{Lean declaration body}
\begin{ReviewVerbatim}
type:
{α : Type u_1} → [Fintype α] → PMF α → (α → ℝ) → ℝ

value:
fun {α : Type u_1} [Fintype α] (μ : PMF α) (f : α → ℝ) => ∑ a : α, (μ a).toReal * f a
\end{ReviewVerbatim}


\hypertarget{governing-declaration-9c38d53b9d664221}{}
\subsection*{\small\ttfamily\raggedright Applied\allowbreak{}Modeling\allowbreak{}Lib.\allowbreak{}pmf\allowbreak{}Prob}
\textbf{Declaration location:} reusable library
\paragraph{Lean declaration body}
\begin{ReviewVerbatim}
type:
{α : Type u_1} → [Fintype α] → PMF α → (p : α → Prop) → [DecidablePred p] → ℝ

value:
fun {α : Type u_1} [Fintype α] (μ : PMF α) (p : α → Prop) [DecidablePred p] =>
  AppliedModelingLib.pmfExp μ fun (a : α) => if p a then 1 else 0
\end{ReviewVerbatim}

\paragraph{Direct retained dependencies}
\begin{itemize}\raggedright
\item \hyperlink{governing-declaration-8e4b4389a4f6454b}{Applied\allowbreak{}Modeling\allowbreak{}Lib.\allowbreak{}pmf\allowbreak{}Exp}
\end{itemize}
\hypertarget{governing-declaration-48f10f2715c04df3}{}
\subsection*{\small\ttfamily\raggedright Applied\allowbreak{}Modeling\allowbreak{}Lib.\allowbreak{}uniform\allowbreak{}PMF}
\textbf{Declaration location:} reusable library
\paragraph{Lean declaration body}
\begin{ReviewVerbatim}
type:
(α : Type u_1) → [Fintype α] → [Nonempty α] → PMF α

value:
fun (α : Type u_1) [Fintype α] [Nonempty α] =>
  PMF.ofFintype (fun (x : α) => (↑(Fintype.card α))⁻¹) (AppliedModelingLib.uniformPMF._proof_1 α)
\end{ReviewVerbatim}


\hypertarget{governing-declaration-6ffc06082b09b7ad}{}
\subsection*{\small\ttfamily\raggedright DSWG24\allowbreak{}Discretization\allowbreak{}Bias.\allowbreak{}Finite.\allowbreak{}feature\allowbreak{}Marginal}
\textbf{Declaration location:} paper-local
\paragraph{Lean declaration body}
\begin{ReviewVerbatim}
type:
{X : Type u_3} → {Y : Type u_4} → [Fintype X] → [DecidableEq X] → [Fintype Y] → [DecidableEq Y] → PMF (X × Y) → PMF X

value:
fun {X : Type u_3} {Y : Type u_4} [Fintype X] [DecidableEq X] [Fintype Y] [DecidableEq Y] (μ : PMF (X × Y)) =>
  PMF.map Prod.fst μ
\end{ReviewVerbatim}


\hypertarget{governing-declaration-6b30364bc30e2f1c}{}
\subsection*{\small\ttfamily\raggedright DSWG24\allowbreak{}Discretization\allowbreak{}Bias.\allowbreak{}Finite.\allowbreak{}label\allowbreak{}Marginal}
\textbf{Declaration location:} paper-local
\paragraph{Lean declaration body}
\begin{ReviewVerbatim}
type:
{X : Type u_3} → {Y : Type u_4} → [Fintype X] → [DecidableEq X] → [Fintype Y] → [DecidableEq Y] → PMF (X × Y) → PMF Y

value:
fun {X : Type u_3} {Y : Type u_4} [Fintype X] [DecidableEq X] [Fintype Y] [DecidableEq Y] (μ : PMF (X × Y)) =>
  PMF.map Prod.snd μ
\end{ReviewVerbatim}


\hypertarget{governing-declaration-7bdd5ebec3f67e16}{}
\subsection*{\small\ttfamily\raggedright DSWG24\allowbreak{}Discretization\allowbreak{}Bias.\allowbreak{}Finite.\allowbreak{}paper\allowbreak{}Aggregate\allowbreak{}Posterior}
\textbf{Declaration location:} paper-local
\paragraph{Lean declaration body}
\begin{ReviewVerbatim}
type:
{X : Type u_3} →
  {Y : Type u_4} → [Fintype X] → [DecidableEq X] → [Fintype Y] → [DecidableEq Y] → PMF (X × Y) → (X → Y → ℝ) → Y → ℝ

value:
fun {X : Type u_3} {Y : Type u_4} [Fintype X] [DecidableEq X] [Fintype Y] [DecidableEq Y] (μ : PMF (X × Y))
    (q : X → Y → ℝ) (y : Y) =>
  AppliedModelingLib.pmfExp (DSWG24DiscretizationBias.Finite.featureMarginal μ) fun (x : X) => q x y
\end{ReviewVerbatim}

\paragraph{Direct retained dependencies}
\begin{itemize}\raggedright
\item \hyperlink{governing-declaration-8e4b4389a4f6454b}{Applied\allowbreak{}Modeling\allowbreak{}Lib.\allowbreak{}pmf\allowbreak{}Exp}
\item \hyperlink{governing-declaration-6ffc06082b09b7ad}{DSWG24\allowbreak{}Discretization\allowbreak{}Bias.\allowbreak{}Finite.\allowbreak{}feature\allowbreak{}Marginal}
\end{itemize}
\hypertarget{governing-declaration-8f84d0f76e66fc93}{}
\subsection*{\small\ttfamily\raggedright DSWG24\allowbreak{}Discretization\allowbreak{}Bias.\allowbreak{}Finite.\allowbreak{}paper\allowbreak{}Marginal\allowbreak{}Label\allowbreak{}Share}
\textbf{Declaration location:} paper-local
\paragraph{Lean declaration body}
\begin{ReviewVerbatim}
type:
{X : Type u_3} →
  {Y : Type u_4} → [Fintype X] → [DecidableEq X] → [Fintype Y] → [DecidableEq Y] → PMF (X × Y) → (X → Y) → Y → ℝ

value:
fun {X : Type u_3} {Y : Type u_4} [Fintype X] [DecidableEq X] [Fintype Y] [DecidableEq Y] (μ : PMF (X × Y))
    (rule : X → Y) (y : Y) =>
  AppliedModelingLib.pmfProb (DSWG24DiscretizationBias.Finite.featureMarginal μ) fun (x : X) => rule x = y
\end{ReviewVerbatim}

\paragraph{Direct retained dependencies}
\begin{itemize}\raggedright
\item \hyperlink{governing-declaration-9c38d53b9d664221}{Applied\allowbreak{}Modeling\allowbreak{}Lib.\allowbreak{}pmf\allowbreak{}Prob}
\item \hyperlink{governing-declaration-6ffc06082b09b7ad}{DSWG24\allowbreak{}Discretization\allowbreak{}Bias.\allowbreak{}Finite.\allowbreak{}feature\allowbreak{}Marginal}
\end{itemize}
\hypertarget{governing-declaration-5d2200dbd9b4b93d}{}
\subsection*{\small\ttfamily\raggedright DSWG24\allowbreak{}Discretization\allowbreak{}Bias.\allowbreak{}Finite.\allowbreak{}paper\allowbreak{}Bias}
\textbf{Declaration location:} paper-local
\paragraph{Lean declaration body}
\begin{ReviewVerbatim}
type:
{X : Type u_3} →
  {Y : Type u_4} →
    [Fintype X] → [DecidableEq X] → [Fintype Y] → [DecidableEq Y] → PMF (X × Y) → (X → Y) → (Y → ℝ) → Y → ℝ

value:
fun {X : Type u_3} {Y : Type u_4} [Fintype X] [DecidableEq X] [Fintype Y] [DecidableEq Y] (μ : PMF (X × Y))
    (rule : X → Y) (pref : Y → ℝ) (y : Y) =>
  DSWG24DiscretizationBias.Finite.paperMarginalLabelShare μ rule y - pref y
\end{ReviewVerbatim}

\paragraph{Direct retained dependencies}
\begin{itemize}\raggedright
\item \hyperlink{governing-declaration-8f84d0f76e66fc93}{DSWG24\allowbreak{}Discretization\allowbreak{}Bias.\allowbreak{}Finite.\allowbreak{}paper\allowbreak{}Marginal\allowbreak{}Label\allowbreak{}Share}
\end{itemize}
\hypertarget{governing-declaration-f3bd0c0016772e11}{}
\subsection*{\small\ttfamily\raggedright DSWG24\allowbreak{}Discretization\allowbreak{}Bias.\allowbreak{}Finite.\allowbreak{}paper\allowbreak{}Aggregate\allowbreak{}Bias}
\textbf{Declaration location:} paper-local
\paragraph{Lean declaration body}
\begin{ReviewVerbatim}
type:
{X : Type u_3} →
  {Y : Type u_4} →
    [Fintype X] → [DecidableEq X] → [Fintype Y] → [DecidableEq Y] → PMF (X × Y) → (X → Y → ℝ) → (X → Y) → Y → ℝ

value:
fun {X : Type u_3} {Y : Type u_4} [Fintype X] [DecidableEq X] [Fintype Y] [DecidableEq Y] (μ : PMF (X × Y))
    (q : X → Y → ℝ) (rule : X → Y) (y : Y) =>
  DSWG24DiscretizationBias.Finite.paperBias μ rule (DSWG24DiscretizationBias.Finite.paperAggregatePosterior μ q) y
\end{ReviewVerbatim}

\paragraph{Direct retained dependencies}
\begin{itemize}\raggedright
\item \hyperlink{governing-declaration-7bdd5ebec3f67e16}{DSWG24\allowbreak{}Discretization\allowbreak{}Bias.\allowbreak{}Finite.\allowbreak{}paper\allowbreak{}Aggregate\allowbreak{}Posterior}
\item \hyperlink{governing-declaration-5d2200dbd9b4b93d}{DSWG24\allowbreak{}Discretization\allowbreak{}Bias.\allowbreak{}Finite.\allowbreak{}paper\allowbreak{}Bias}
\end{itemize}
\hypertarget{governing-declaration-9de2c431ac58b2a3}{}
\subsection*{\small\ttfamily\raggedright DSWG24\allowbreak{}Discretization\allowbreak{}Bias.\allowbreak{}Finite.\allowbreak{}paper\allowbreak{}Theorem1\allowbreak{}Tight\allowbreak{}Argmax}
\textbf{Declaration location:} paper-local
\paragraph{Lean declaration body}
\begin{ReviewVerbatim}
type:
Unit → Fin 2

value:
fun (x : Unit) => 0
\end{ReviewVerbatim}


\hypertarget{governing-declaration-d4243b460c2c568f}{}
\subsection*{\small\ttfamily\raggedright DSWG24\allowbreak{}Discretization\allowbreak{}Bias.\allowbreak{}Finite.\allowbreak{}paper\allowbreak{}Theorem1\allowbreak{}Tight\allowbreak{}Mu}
\textbf{Declaration location:} paper-local
\paragraph{Lean declaration body}
\begin{ReviewVerbatim}
type:
PMF (Unit × Fin 2)

value:
AppliedModelingLib.uniformPMF (Unit × Fin 2)
\end{ReviewVerbatim}

\paragraph{Direct retained dependencies}
\begin{itemize}\raggedright
\item \hyperlink{governing-declaration-48f10f2715c04df3}{Applied\allowbreak{}Modeling\allowbreak{}Lib.\allowbreak{}uniform\allowbreak{}PMF}
\end{itemize}
\hypertarget{governing-declaration-193c98b93e9de012}{}
\subsection*{\small\ttfamily\raggedright DSWG24\allowbreak{}Discretization\allowbreak{}Bias.\allowbreak{}Finite.\allowbreak{}paper\allowbreak{}Theorem1\allowbreak{}Tight\allowbreak{}Q}
\textbf{Declaration location:} paper-local
\paragraph{Lean declaration body}
\begin{ReviewVerbatim}
type:
Unit → Fin 2 → ℝ

value:
fun (x : Unit) (x_1 : Fin 2) => 1 / 2
\end{ReviewVerbatim}


\hypertarget{governing-declaration-41ac4bb8a2edc19c}{}
\subsection*{\small\ttfamily\raggedright DSWG24\allowbreak{}Discretization\allowbreak{}Bias.\allowbreak{}Pareto.\allowbreak{}Pareto\allowbreak{}Dominates}
\textbf{Declaration location:} paper-local
\paragraph{Lean declaration body}
\begin{ReviewVerbatim}
type:
{Rule : Type u_1} → (Rule → ℝ) → (Rule → ℝ) → Rule → Rule → Prop

value:
fun {Rule : Type u_1} (accuracy fidelity : Rule → ℝ) (better worse : Rule) =>
  accuracy worse ≤ accuracy better ∧
    fidelity worse ≤ fidelity better ∧ (accuracy worse < accuracy better ∨ fidelity worse < fidelity better)
\end{ReviewVerbatim}


\hypertarget{governing-declaration-90962900f461a9f9}{}
\subsection*{\small\ttfamily\raggedright DSWG24\allowbreak{}Discretization\allowbreak{}Bias.\allowbreak{}Pareto.\allowbreak{}Pareto\allowbreak{}Optimal}
\textbf{Declaration location:} paper-local
\paragraph{Lean declaration body}
\begin{ReviewVerbatim}
type:
{Rule : Type u_1} → (Rule → ℝ) → (Rule → ℝ) → Rule → Prop

value:
fun {Rule : Type u_1} (accuracy fidelity : Rule → ℝ) (rule : Rule) =>
  ¬∃ (better : Rule), DSWG24DiscretizationBias.Pareto.ParetoDominates accuracy fidelity better rule
\end{ReviewVerbatim}

\paragraph{Direct retained dependencies}
\begin{itemize}\raggedright
\item \hyperlink{governing-declaration-41ac4bb8a2edc19c}{DSWG24\allowbreak{}Discretization\allowbreak{}Bias.\allowbreak{}Pareto.\allowbreak{}Pareto\allowbreak{}Dominates}
\end{itemize}
\hypertarget{governing-declaration-d0facd526e2aa2ff}{}
\subsection*{\small\ttfamily\raggedright DSWG24\allowbreak{}Discretization\allowbreak{}Bias.\allowbreak{}Pareto.\allowbreak{}weighted\allowbreak{}Objective}
\textbf{Declaration location:} paper-local
\paragraph{Lean declaration body}
\begin{ReviewVerbatim}
type:
{Rule : Type u_1} → ℝ → (Rule → ℝ) → (Rule → ℝ) → Rule → ℝ

value:
fun {Rule : Type u_1} (γ : ℝ) (accuracy fidelity : Rule → ℝ) (rule : Rule) =>
  γ * accuracy rule + (1 - γ) * fidelity rule
\end{ReviewVerbatim}


\hypertarget{governing-declaration-71a4a73d8c8948eb}{}
\subsection*{\small\ttfamily\raggedright DSWG24\allowbreak{}Discretization\allowbreak{}Bias.\allowbreak{}Proof\allowbreak{}Bridge.\allowbreak{}aggregate\allowbreak{}Posterior}
\textbf{Declaration location:} paper-local
\paragraph{Lean declaration body}
\begin{ReviewVerbatim}
type:
{N K : ℕ} → (Fin N → Fin K → ℝ) → Fin K → ℝ

value:
fun {N K : ℕ} (q : Fin N → Fin K → ℝ) (y : Fin K) => (∑ i : Fin N, q i y) / ↑N
\end{ReviewVerbatim}


\hypertarget{governing-declaration-2e16128ba60f3c18}{}
\subsection*{\small\ttfamily\raggedright DSWG24\allowbreak{}Discretization\allowbreak{}Bias.\allowbreak{}Proof\allowbreak{}Bridge.\allowbreak{}classifier\allowbreak{}MAE}
\textbf{Declaration location:} paper-local
\paragraph{Lean declaration body}
\begin{ReviewVerbatim}
type:
{X : Type u_1} →
  {Y : Type u_2} → [Fintype X] → [DecidableEq X] → [Fintype Y] → [DecidableEq Y] → PMF (X × Y) → (X → Y → ℝ) → ℝ

value:
fun {X : Type u_1} {Y : Type u_2} [Fintype X] [DecidableEq X] [Fintype Y] [DecidableEq Y] (μ : PMF (X × Y))
    (q : X → Y → ℝ) =>
  AppliedModelingLib.pmfExp (DSWG24DiscretizationBias.Finite.featureMarginal μ) fun (x : X) =>
    ∑ y : Y, q x y * (1 - q x y)
\end{ReviewVerbatim}

\paragraph{Direct retained dependencies}
\begin{itemize}\raggedright
\item \hyperlink{governing-declaration-8e4b4389a4f6454b}{Applied\allowbreak{}Modeling\allowbreak{}Lib.\allowbreak{}pmf\allowbreak{}Exp}
\item \hyperlink{governing-declaration-6ffc06082b09b7ad}{DSWG24\allowbreak{}Discretization\allowbreak{}Bias.\allowbreak{}Finite.\allowbreak{}feature\allowbreak{}Marginal}
\end{itemize}
\hypertarget{governing-declaration-d2886661b52542f6}{}
\subsection*{\small\ttfamily\raggedright DSWG24\allowbreak{}Discretization\allowbreak{}Bias.\allowbreak{}Proof\allowbreak{}Bridge.\allowbreak{}continuous\allowbreak{}Joint\allowbreak{}Aggregate\allowbreak{}Posterior}
\textbf{Declaration location:} paper-local
\paragraph{Lean declaration body}
\begin{ReviewVerbatim}
type:
{X : Type u_1} → {Y : Type u_2} → [inst : MeasurableSpace (X × Y)] → MeasureTheory.Measure (X × Y) → (X → Y → ℝ) → Y → ℝ

value:
fun {X : Type u_1} {Y : Type u_2} [MeasurableSpace (X × Y)] (μ : MeasureTheory.Measure (X × Y)) (q : X → Y → ℝ)
    (y : Y) =>
  ∫ (xy : X × Y), q xy.1 y ∂μ
\end{ReviewVerbatim}


\hypertarget{governing-declaration-302610b6fa0b99db}{}
\subsection*{\small\ttfamily\raggedright DSWG24\allowbreak{}Discretization\allowbreak{}Bias.\allowbreak{}Proof\allowbreak{}Bridge.\allowbreak{}continuous\allowbreak{}Joint\allowbreak{}Aggregate\allowbreak{}Bias}
\textbf{Declaration location:} paper-local
\paragraph{Lean declaration body}
\begin{ReviewVerbatim}
type:
{X : Type u_1} →
  {Y : Type u_2} →
    [inst : MeasurableSpace (X × Y)] → [DecidableEq Y] → MeasureTheory.Measure (X × Y) → (X → Y → ℝ) → (X → Y) → Y → ℝ

value:
fun {X : Type u_1} {Y : Type u_2} [MeasurableSpace (X × Y)] [DecidableEq Y] (μ : MeasureTheory.Measure (X × Y))
    (q : X → Y → ℝ) (rule : X → Y) (y : Y) =>
  ∫ (xy : X × Y), if rule xy.1 = y then 1 else 0 ∂μ -
    DSWG24DiscretizationBias.ProofBridge.continuousJointAggregatePosterior μ q y
\end{ReviewVerbatim}

\paragraph{Direct retained dependencies}
\begin{itemize}\raggedright
\item \hyperlink{governing-declaration-d2886661b52542f6}{DSWG24\allowbreak{}Discretization\allowbreak{}Bias.\allowbreak{}Proof\allowbreak{}Bridge.\allowbreak{}continuous\allowbreak{}Joint\allowbreak{}Aggregate\allowbreak{}Posterior}
\end{itemize}
\hypertarget{governing-declaration-d07991e4ae99da3b}{}
\subsection*{\small\ttfamily\raggedright DSWG24\allowbreak{}Discretization\allowbreak{}Bias.\allowbreak{}Proof\allowbreak{}Bridge.\allowbreak{}continuous\allowbreak{}Joint\allowbreak{}Classifier\allowbreak{}MAE}
\textbf{Declaration location:} paper-local
\paragraph{Lean declaration body}
\begin{ReviewVerbatim}
type:
{X : Type u_1} →
  {Y : Type u_2} → [inst : MeasurableSpace (X × Y)] → [Fintype Y] → MeasureTheory.Measure (X × Y) → (X → Y → ℝ) → ℝ

value:
fun {X : Type u_1} {Y : Type u_2} [MeasurableSpace (X × Y)] [Fintype Y] (μ : MeasureTheory.Measure (X × Y))
    (q : X → Y → ℝ) =>
  ∫ (xy : X × Y), ∑ y : Y, q xy.1 y * (1 - q xy.1 y) ∂μ
\end{ReviewVerbatim}


\hypertarget{governing-declaration-7df1fb937c21cdf5}{}
\subsection*{\small\ttfamily\raggedright DSWG24\allowbreak{}Discretization\allowbreak{}Bias.\allowbreak{}Proof\allowbreak{}Bridge.\allowbreak{}continuous\allowbreak{}Joint\allowbreak{}Prior\allowbreak{}Bias}
\textbf{Declaration location:} paper-local
\paragraph{Lean declaration body}
\begin{ReviewVerbatim}
type:
{X : Type u_1} →
  {Y : Type u_2} → [inst : MeasurableSpace (X × Y)] → [DecidableEq Y] → MeasureTheory.Measure (X × Y) → (X → Y) → Y → ℝ

value:
fun {X : Type u_1} {Y : Type u_2} [MeasurableSpace (X × Y)] [DecidableEq Y] (μ : MeasureTheory.Measure (X × Y))
    (rule : X → Y) (y : Y) =>
  ∫ (xy : X × Y), if rule xy.1 = y then 1 else 0 ∂μ - ∫ (xy : X × Y), if xy.2 = y then 1 else 0 ∂μ
\end{ReviewVerbatim}


\hypertarget{governing-declaration-8e0892ece8e4e41c}{}
\subsection*{\small\ttfamily\raggedright DSWG24\allowbreak{}Discretization\allowbreak{}Bias.\allowbreak{}Proof\allowbreak{}Bridge.\allowbreak{}expected\allowbreak{}Objective}
\textbf{Declaration location:} paper-local
\paragraph{Lean declaration body}
\begin{ReviewVerbatim}
type:
{ω : Type u_1} →
  {σ : Type u_2} →
    {N K : ℕ} →
      [NeZero K] →
        ((ω → ℝ) → ℝ) → (ω → σ) → (ω → Fin N → Fin K) → ℝ → (σ → (Fin N → Fin K) → ℝ) → (σ → Fin N → Fin K) → ℝ

value:
fun {ω : Type u_1} {σ : Type u_2} {N K : ℕ} [NeZero K] (expect : (ω → ℝ) → ℝ) (observedDataset : ω → σ)
    (trueLabels : ω → Fin N → Fin K) (γ : ℝ) (fidelityTerm : σ → (Fin N → Fin K) → ℝ) (rule : σ → Fin N → Fin K) =>
  γ * AppliedModelingLib.Decision.expectedDecisionAccuracy expect observedDataset trueLabels rule +
    (1 - γ) * AppliedModelingLib.Decision.expectedObjective expect observedDataset fidelityTerm rule
\end{ReviewVerbatim}

\paragraph{Direct retained dependencies}
\begin{itemize}\raggedright
\item \hyperlink{governing-declaration-61110afe97ed1775}{Applied\allowbreak{}Modeling\allowbreak{}Lib.\allowbreak{}Decision.\allowbreak{}expected\allowbreak{}Decision\allowbreak{}Accuracy}
\item \hyperlink{governing-declaration-71856cc0344b281a}{Applied\allowbreak{}Modeling\allowbreak{}Lib.\allowbreak{}Decision.\allowbreak{}expected\allowbreak{}Objective}
\end{itemize}
\hypertarget{governing-declaration-af44c4b45082ddcb}{}
\subsection*{\small\ttfamily\raggedright DSWG24\allowbreak{}Discretization\allowbreak{}Bias.\allowbreak{}Proof\allowbreak{}Bridge.\allowbreak{}is\allowbreak{}First\allowbreak{}Argmax}
\textbf{Declaration location:} paper-local
\paragraph{Lean declaration body}
\begin{ReviewVerbatim}
type:
{K : ℕ} → (Fin K → ℝ) → Fin K → Prop

value:
fun {K : ℕ} (score : Fin K → ℝ) (chosen : Fin K) =>
  (∀ (y : Fin K), score y ≤ score chosen) ∧ ∀ (y : Fin K), score y = score chosen → ↑chosen ≤ ↑y
\end{ReviewVerbatim}


\hypertarget{governing-declaration-105d567a51316da3}{}
\subsection*{\small\ttfamily\raggedright DSWG24\allowbreak{}Discretization\allowbreak{}Bias.\allowbreak{}Proof\allowbreak{}Bridge.\allowbreak{}marginal\allowbreak{}Label\allowbreak{}Share}
\textbf{Declaration location:} paper-local
\paragraph{Lean declaration body}
\begin{ReviewVerbatim}
type:
{N K : ℕ} → (Fin N → Fin K) → Fin K → ℝ

value:
fun {N K : ℕ} (decision : Fin N → Fin K) (y : Fin K) => (∑ i : Fin N, if decision i = y then 1 else 0) / ↑N
\end{ReviewVerbatim}


\hypertarget{governing-declaration-63f8cdbf3ab80aed}{}
\subsection*{\small\ttfamily\raggedright DSWG24\allowbreak{}Discretization\allowbreak{}Bias.\allowbreak{}Proof\allowbreak{}Bridge.\allowbreak{}bias}
\textbf{Declaration location:} paper-local
\paragraph{Lean declaration body}
\begin{ReviewVerbatim}
type:
{N K : ℕ} → (Fin N → Fin K) → (Fin K → ℝ) → Fin K → ℝ

value:
fun {N K : ℕ} (decision : Fin N → Fin K) (pref : Fin K → ℝ) (y : Fin K) =>
  DSWG24DiscretizationBias.ProofBridge.marginalLabelShare decision y - pref y
\end{ReviewVerbatim}

\paragraph{Direct retained dependencies}
\begin{itemize}\raggedright
\item \hyperlink{governing-declaration-105d567a51316da3}{DSWG24\allowbreak{}Discretization\allowbreak{}Bias.\allowbreak{}Proof\allowbreak{}Bridge.\allowbreak{}marginal\allowbreak{}Label\allowbreak{}Share}
\end{itemize}
\hypertarget{governing-declaration-2b57ec17c49d7147}{}
\subsection*{\small\ttfamily\raggedright DSWG24\allowbreak{}Discretization\allowbreak{}Bias.\allowbreak{}Proof\allowbreak{}Bridge.\allowbreak{}fidelity}
\textbf{Declaration location:} paper-local
\paragraph{Lean declaration body}
\begin{ReviewVerbatim}
type:
{N K : ℕ} → (Fin N → Fin K) → (Fin K → ℝ) → ℝ

value:
fun {N K : ℕ} (decision : Fin N → Fin K) (pref : Fin K → ℝ) =>
  -∑ y : Fin K, |DSWG24DiscretizationBias.ProofBridge.bias decision pref y|
\end{ReviewVerbatim}

\paragraph{Direct retained dependencies}
\begin{itemize}\raggedright
\item \hyperlink{governing-declaration-63f8cdbf3ab80aed}{DSWG24\allowbreak{}Discretization\allowbreak{}Bias.\allowbreak{}Proof\allowbreak{}Bridge.\allowbreak{}bias}
\end{itemize}
\hypertarget{governing-declaration-78f0292bcb487fbc}{}
\subsection*{\small\ttfamily\raggedright DSWG24\allowbreak{}Discretization\allowbreak{}Bias.\allowbreak{}Proof\allowbreak{}Bridge.\allowbreak{}equation1\allowbreak{}Objective}
\textbf{Declaration location:} paper-local
\paragraph{Lean declaration body}
\begin{ReviewVerbatim}
type:
{N K : ℕ} → ℝ → (Fin N → Fin K → ℝ) → (Fin K → ℝ) → (Fin N → Fin K) → ℝ

value:
fun {N K : ℕ} (γ : ℝ) (q : Fin N → Fin K → ℝ) (pref : Fin K → ℝ) (decision : Fin N → Fin K) =>
  γ * AppliedModelingLib.Decision.averageScore q decision +
    (1 - γ) * DSWG24DiscretizationBias.ProofBridge.fidelity decision pref
\end{ReviewVerbatim}

\paragraph{Direct retained dependencies}
\begin{itemize}\raggedright
\item \hyperlink{governing-declaration-3088d80cd010ba5f}{Applied\allowbreak{}Modeling\allowbreak{}Lib.\allowbreak{}Decision.\allowbreak{}average\allowbreak{}Score}
\item \hyperlink{governing-declaration-2b57ec17c49d7147}{DSWG24\allowbreak{}Discretization\allowbreak{}Bias.\allowbreak{}Proof\allowbreak{}Bridge.\allowbreak{}fidelity}
\end{itemize}
\hypertarget{governing-declaration-28e1fcee264c130f}{}
\subsection*{\small\ttfamily\raggedright DSWG24\allowbreak{}Discretization\allowbreak{}Bias.\allowbreak{}Proof\allowbreak{}Bridge.\allowbreak{}perfect\allowbreak{}Classifier\allowbreak{}On\allowbreak{}Support}
\textbf{Declaration location:} paper-local
\paragraph{Lean declaration body}
\begin{ReviewVerbatim}
type:
{X : Type u_1} →
  {Y : Type u_2} →
    [Fintype X] → [DecidableEq X] → [Fintype Y] → [DecidableEq Y] → PMF (X × Y) → (X → Y → ℝ) → (X → Y) → Prop

value:
fun {X : Type u_1} {Y : Type u_2} [Fintype X] [DecidableEq X] [Fintype Y] [DecidableEq Y] (μ : PMF (X × Y))
    (q : X → Y → ℝ) (rule : X → Y) =>
  ∀ (x : X) (y : Y), μ (x, y) ≠ 0 → rule x = y ∧ q x y = 1 ∧ ∀ (a : Y), a ≠ y → q x a = 0
\end{ReviewVerbatim}


\hypertarget{governing-declaration-8260e4a6984dc105}{}
\subsection*{\small\ttfamily\raggedright DSWG24\allowbreak{}Discretization\allowbreak{}Bias.\allowbreak{}Proof\allowbreak{}Bridge.\allowbreak{}prior}
\textbf{Declaration location:} paper-local
\paragraph{Lean declaration body}
\begin{ReviewVerbatim}
type:
{X : Type u_1} → {Y : Type u_2} → [Fintype X] → [DecidableEq X] → [Fintype Y] → [DecidableEq Y] → PMF (X × Y) → Y → ℝ

value:
fun {X : Type u_1} {Y : Type u_2} [Fintype X] [DecidableEq X] [Fintype Y] [DecidableEq Y] (μ : PMF (X × Y)) (y : Y) =>
  ((DSWG24DiscretizationBias.Finite.labelMarginal μ) y).toReal
\end{ReviewVerbatim}

\paragraph{Direct retained dependencies}
\begin{itemize}\raggedright
\item \hyperlink{governing-declaration-6b30364bc30e2f1c}{DSWG24\allowbreak{}Discretization\allowbreak{}Bias.\allowbreak{}Finite.\allowbreak{}label\allowbreak{}Marginal}
\end{itemize}
\hypertarget{governing-declaration-9631b3d6800a2bd3}{}
\subsection*{\small\ttfamily\raggedright DSWG24\allowbreak{}Discretization\allowbreak{}Bias.\allowbreak{}Proof\allowbreak{}Bridge.\allowbreak{}source\allowbreak{}P\allowbreak{}Nq}
\textbf{Declaration location:} paper-local
\paragraph{Lean declaration body}
\begin{ReviewVerbatim}
type:
{Sample : Type u_1} → {N K : ℕ} → (Sample → Fin N → Fin K → ℝ) → (Sample → Fin K) → (Sample → Fin K → ℝ) → Prop

value:
fun {Sample : Type u_1} {N K : ℕ} (posterior : Sample → Fin N → Fin K → ℝ) (aggregateArgmax : Sample → Fin K)
    (prefAt : Sample → Fin K → ℝ) =>
  DSWG24DiscretizationBias.ProofBridge.referenceSimplexAt prefAt ∧
    (∀ (sample : Sample),
        DSWG24DiscretizationBias.ProofBridge.isAggregateFirstArgmax (posterior sample) (aggregateArgmax sample)) ∧
      ∀ (sample : Sample), 1 / ↑N ≤ prefAt sample (aggregateArgmax sample)
\end{ReviewVerbatim}

\paragraph{Direct retained dependencies}
\begin{itemize}\raggedright
\item \hyperlink{paper-prerequisite-46e579d6fcd53814}{DSWG24\allowbreak{}Discretization\allowbreak{}Bias.\allowbreak{}Proof\allowbreak{}Bridge.\allowbreak{}is\allowbreak{}Aggregate\allowbreak{}First\allowbreak{}Argmax}
\item \hyperlink{paper-prerequisite-e9ff5dcccf634407}{DSWG24\allowbreak{}Discretization\allowbreak{}Bias.\allowbreak{}Proof\allowbreak{}Bridge.\allowbreak{}reference\allowbreak{}Simplex\allowbreak{}At}
\end{itemize}
\hypertarget{governing-declaration-fd6e69b98a707ae7}{}
\subsection*{\small\ttfamily\raggedright DSWG24\allowbreak{}Discretization\allowbreak{}Bias.\allowbreak{}Theorem2iii.\allowbreak{}dataset\allowbreak{}Accuracy\allowbreak{}Score}
\textbf{Declaration location:} paper-local
\paragraph{Lean declaration body}
\begin{ReviewVerbatim}
type:
{N K : ℕ} → (Fin N → Fin K → ℝ) → (Fin N → Fin K) → ℝ

value:
fun {N K : ℕ} (posterior : Fin N → Fin K → ℝ) (decision : Fin N → Fin K) =>
  AppliedModelingLib.Decision.averageScore posterior decision
\end{ReviewVerbatim}

\paragraph{Direct retained dependencies}
\begin{itemize}\raggedright
\item \hyperlink{governing-declaration-3088d80cd010ba5f}{Applied\allowbreak{}Modeling\allowbreak{}Lib.\allowbreak{}Decision.\allowbreak{}average\allowbreak{}Score}
\end{itemize}
\hypertarget{governing-declaration-81ede32916dad99e}{}
\subsection*{\small\ttfamily\raggedright DSWG24\allowbreak{}Discretization\allowbreak{}Bias.\allowbreak{}Theorem2iii.\allowbreak{}dataset\allowbreak{}Label\allowbreak{}Share}
\textbf{Declaration location:} paper-local
\paragraph{Lean declaration body}
\begin{ReviewVerbatim}
type:
{N K : ℕ} → (Fin N → Fin K) → Fin K → ℝ

value:
fun {N K : ℕ} (decision : Fin N → Fin K) (y : Fin K) => (∑ i : Fin N, if decision i = y then 1 else 0) / ↑N
\end{ReviewVerbatim}


\hypertarget{governing-declaration-2a64863d3ac17c16}{}
\subsection*{\small\ttfamily\raggedright DSWG24\allowbreak{}Discretization\allowbreak{}Bias.\allowbreak{}Theorem2iii.\allowbreak{}dataset\allowbreak{}Fidelity}
\textbf{Declaration location:} paper-local
\paragraph{Lean declaration body}
\begin{ReviewVerbatim}
type:
{N K : ℕ} → (Fin K → ℝ) → (Fin N → Fin K) → ℝ

value:
fun {N K : ℕ} (pref : Fin K → ℝ) (decision : Fin N → Fin K) =>
  -∑ y : Fin K, |DSWG24DiscretizationBias.Theorem2iii.datasetLabelShare decision y - pref y|
\end{ReviewVerbatim}

\paragraph{Direct retained dependencies}
\begin{itemize}\raggedright
\item \hyperlink{governing-declaration-81ede32916dad99e}{DSWG24\allowbreak{}Discretization\allowbreak{}Bias.\allowbreak{}Theorem2iii.\allowbreak{}dataset\allowbreak{}Label\allowbreak{}Share}
\end{itemize}
\hypertarget{governing-declaration-daac633422a44729}{}
\subsection*{\small\ttfamily\raggedright DSWG24\allowbreak{}Discretization\allowbreak{}Bias.\allowbreak{}Theorem2iii.\allowbreak{}expected\allowbreak{}Dataset\allowbreak{}Accuracy}
\textbf{Declaration location:} paper-local
\paragraph{Lean declaration body}
\begin{ReviewVerbatim}
type:
{Sample : Type u_1} →
  [Fintype Sample] →
    [DecidableEq Sample] → {N K : ℕ} → PMF Sample → (Sample → Fin N → Fin K → ℝ) → (Sample → Fin N → Fin K) → ℝ

value:
fun {Sample : Type u_1} [Fintype Sample] [DecidableEq Sample] {N K : ℕ} (sampleLaw : PMF Sample)
    (posterior : Sample → Fin N → Fin K → ℝ) (policy : Sample → Fin N → Fin K) =>
  AppliedModelingLib.pmfExp sampleLaw fun (sample : Sample) =>
    DSWG24DiscretizationBias.Theorem2iii.datasetAccuracyScore (posterior sample) (policy sample)
\end{ReviewVerbatim}

\paragraph{Direct retained dependencies}
\begin{itemize}\raggedright
\item \hyperlink{governing-declaration-8e4b4389a4f6454b}{Applied\allowbreak{}Modeling\allowbreak{}Lib.\allowbreak{}pmf\allowbreak{}Exp}
\item \hyperlink{governing-declaration-fd6e69b98a707ae7}{DSWG24\allowbreak{}Discretization\allowbreak{}Bias.\allowbreak{}Theorem2iii.\allowbreak{}dataset\allowbreak{}Accuracy\allowbreak{}Score}
\end{itemize}
\hypertarget{governing-declaration-8158479406f1d885}{}
\subsection*{\small\ttfamily\raggedright DSWG24\allowbreak{}Discretization\allowbreak{}Bias.\allowbreak{}Theorem2iii.\allowbreak{}expected\allowbreak{}Dataset\allowbreak{}Fidelity\allowbreak{}At}
\textbf{Declaration location:} paper-local
\paragraph{Lean declaration body}
\begin{ReviewVerbatim}
type:
{Sample : Type u_1} →
  [Fintype Sample] → [DecidableEq Sample] → {N K : ℕ} → PMF Sample → (Sample → Fin K → ℝ) → (Sample → Fin N → Fin K) → ℝ

value:
fun {Sample : Type u_1} [Fintype Sample] [DecidableEq Sample] {N K : ℕ} (sampleLaw : PMF Sample)
    (prefAt : Sample → Fin K → ℝ) (policy : Sample → Fin N → Fin K) =>
  AppliedModelingLib.pmfExp sampleLaw fun (sample : Sample) =>
    DSWG24DiscretizationBias.Theorem2iii.datasetFidelity (prefAt sample) (policy sample)
\end{ReviewVerbatim}

\paragraph{Direct retained dependencies}
\begin{itemize}\raggedright
\item \hyperlink{governing-declaration-8e4b4389a4f6454b}{Applied\allowbreak{}Modeling\allowbreak{}Lib.\allowbreak{}pmf\allowbreak{}Exp}
\item \hyperlink{governing-declaration-2a64863d3ac17c16}{DSWG24\allowbreak{}Discretization\allowbreak{}Bias.\allowbreak{}Theorem2iii.\allowbreak{}dataset\allowbreak{}Fidelity}
\end{itemize}
\hypertarget{governing-declaration-b906fbb2f231ab40}{}
\subsection*{\small\ttfamily\raggedright DSWG24\allowbreak{}Discretization\allowbreak{}Bias.\allowbreak{}Theorem2iii.\allowbreak{}iid\allowbreak{}Sample\allowbreak{}PMF}
\textbf{Declaration location:} paper-local
\paragraph{Lean declaration body}
\begin{ReviewVerbatim}
type:
{X : Type u_1} → [Fintype X] → [DecidableEq X] → PMF X → (N : ℕ) → PMF (Fin N → X)

value:
fun {X : Type u_1} [Fintype X] [DecidableEq X] (μ : PMF X) (N : ℕ) =>
  PMF.ofFintype (fun (sample : Fin N → X) => ∏ i : Fin N, μ (sample i))
    (DSWG24DiscretizationBias.Theorem2iii.iidSamplePMF._proof_1 μ N)
\end{ReviewVerbatim}


\hypertarget{governing-declaration-05d42de5dba2851d}{}
\subsection*{\small\ttfamily\raggedright DSWG24\allowbreak{}Discretization\allowbreak{}Bias.\allowbreak{}Theorem2iii.\allowbreak{}sampled\allowbreak{}Decision}
\textbf{Declaration location:} paper-local
\paragraph{Lean declaration body}
\begin{ReviewVerbatim}
type:
{X : Type u_1} → {N K : ℕ} → (X → Fin K) → (Fin N → X) → Fin N → Fin K

value:
fun {X : Type u_1} {N K : ℕ} (rule : X → Fin K) (sample : Fin N → X) (a : Fin N) => rule (sample a)
\end{ReviewVerbatim}


\hypertarget{governing-declaration-457b4184b0081169}{}
\subsection*{\small\ttfamily\raggedright DSWG24\allowbreak{}Discretization\allowbreak{}Bias.\allowbreak{}Theorem2iii.\allowbreak{}sampled\allowbreak{}Posterior}
\textbf{Declaration location:} paper-local
\paragraph{Lean declaration body}
\begin{ReviewVerbatim}
type:
{X : Type u_1} → {N K : ℕ} → (X → Fin K → ℝ) → (Fin N → X) → Fin N → Fin K → ℝ

value:
fun {X : Type u_1} {N K : ℕ} (posterior : X → Fin K → ℝ) (sample : Fin N → X) (a : Fin N) (a_1 : Fin K) =>
  posterior (sample a) a_1
\end{ReviewVerbatim}


\clearpage
\hypertarget{paper-prerequisite-section-5ef5ef0364b6939c}{}
\section*{Paper-specific semantic prerequisites}
\hypertarget{paper-prerequisite-c793e07e52a48893}{}
\subsection*{\small\ttfamily\raggedright DSWG24\allowbreak{}Discretization\allowbreak{}Bias.\allowbreak{}Proof\allowbreak{}Bridge.\allowbreak{}bayes\allowbreak{}Optimal}
\noindent\textbf{Paper-local declaration:}\par
{\footnotesize\ttfamily\raggedright DSWG24\allowbreak{}Discretization\allowbreak{}Bias.\allowbreak{}Proof\allowbreak{}Bridge.\allowbreak{}bayes\allowbreak{}Optimal\par}
\textbf{Source locator:} {\footnotesize\raggedright cited publication, lines 597-\allowbreak{}607\par}
\paragraph{Verbatim paper-source connection}
\begin{ReviewVerbatim}
3.1.1    Continuous classifier and discretization rules
Suppose we have a machine learning classifier q (e.g., trained on some other dataset) that makes continuous
probability predictions for each class, given the data. Let q(y, x) denote the probability placed on class y



                                                       10

[form-feed in source extraction]
given data x. For example, the Bayes optimal classifier is
                                               q ∗ (y, x) ≜ Pr(Y = y|X = x).
We will refer to p as the prior and q as our learned posterior. Let the simplex of possible probability vectors
over the support Y be △Y , and let q(x) ∈ △Y be the vector of probabilities (q(y, x))y .
\end{ReviewVerbatim}



\paragraph{Lean-expanded paper semantic target}
\begin{ReviewVerbatim}
type:
{Ω : Type u_1} →
  {σ : Type u_2} →
    [Fintype Ω] →
      [DecidableEq Ω] →
        [Fintype σ] →
          [DecidableEq σ] → {N K : ℕ} → PMF Ω → (Ω → σ) → (Ω → Fin N → Fin K) → (σ → Fin N → Fin K → ℝ) → Prop

value:
fun {Ω : Type u_1} {σ : Type u_2} [Fintype Ω] [DecidableEq Ω] [Fintype σ] [DecidableEq σ] {N K : ℕ} (μ : PMF Ω)
    (observedDataset : Ω → σ) (trueLabels : Ω → Fin N → Fin K) (posterior : σ → Fin N → Fin K → ℝ) =>
  ∀ (xs : σ) (i : Fin N) (y : Fin K),
    (AppliedModelingLib.pmfProb μ fun (w : Ω) => observedDataset w = xs ∧ trueLabels w i = y) =
      posterior xs i y * AppliedModelingLib.pmfProb μ fun (w : Ω) => observedDataset w = xs
\end{ReviewVerbatim}

\paragraph{Direct retained dependencies}
\begin{itemize}\raggedright
\item \hyperlink{governing-declaration-9c38d53b9d664221}{Applied\allowbreak{}Modeling\allowbreak{}Lib.\allowbreak{}pmf\allowbreak{}Prob}
\end{itemize}
\paragraph{Recorded source-to-declaration screening}
\textbf{Verdict:} \texttt{matches}
\\
{\raggedright
\textbf{Reason:} The source defines the Bayes classifier by q-\allowbreak{}star(y,x)=Pr(Y=y given X=x);\allowbreak{} the expanded finite target states the equivalent atom factorization Pr(X=xs and Y\_\allowbreak{}i=y)=posterior(xs,i,y) times Pr(X=xs), preserving every positive-\allowbreak{}mass conditional and leaving zero-\allowbreak{}mass conditionals unconstrained.\allowbreak{}
\par}
\reviewmatch{Matches source input}
\noindent\textbf{Reviewer annotation}\par
\reviewerbox
\clearpage
\hypertarget{paper-prerequisite-c94cbc772446ffec}{}
\subsection*{\small\ttfamily\raggedright DSWG24\allowbreak{}Discretization\allowbreak{}Bias.\allowbreak{}Proof\allowbreak{}Bridge.\allowbreak{}calibrated}
\noindent\textbf{Paper-local declaration:}\par
{\footnotesize\ttfamily\raggedright DSWG24\allowbreak{}Discretization\allowbreak{}Bias.\allowbreak{}Proof\allowbreak{}Bridge.\allowbreak{}calibrated\par}
\textbf{Source locator:} {\footnotesize\raggedright cited publication, lines 635-\allowbreak{}658\par}
\paragraph{Verbatim paper-source connection}
\begin{ReviewVerbatim}
3.2     Theoretical Results: Argmax bias and continuous classifier
We first study the relationship between argmax bias and the properties of the continuous classifier. We show
that discretization bias can occur even with unbiased (calibrated) continuous classifiers q, as long as it is
uncertain (has non-zero error).
    The following result holds for all calibrated classifiers. A given classifier q is calibrated when its contin-
uous predictions are correct on average, Pr(Y = y|q(y, x) = c) = c.10 Note that the Bayes optimal classifier
is calibrated.
Theorem 1. Argmax bias depends on predictive uncertainty, i.e., how much information features x provide
about the true class label. Consider calibrated classifier q and the argmax decision rule Dargmax . Let N > K
and consider a reference distribution of either the aggregate posterior or the prior. Then,
  (i). Suppose x provides no information and q(y, x) = Pr(y) for all x, y. Then, amplification bias is
       maximized and all data points are classified as a plurality class:
                                                  (
                                                    1 − Pr(y), if y = arg maxw p(w),
                            BIAS(y, Dargmax ) =
                                                    − Pr(y),      otherwise.
   9 Note that, though related, mean average error of the continuous predictor q(y, x) is not equivalent to expected zero-one

loss of the discrete decisions. The latter, our measure of decision accuracy, also depends on how the continuous predictor is
discretized.                                n                                           o
  10 For all c of non-zero measure: ∀c ∈
                                                 R
                                             c : (x,y) I[q(y, x) = c]f (x, y)d(y, x) > 0 . Further note that calibration implies
Bayesian consistency, i.e., the continuous model cumulatively assigns the proper mass to each class y: EX [q(y, x)] = Pr(y).
\end{ReviewVerbatim}

% report-context-presentation-sha256: dc1fa9810ac3a973e2d1a457fa41cd1dc4334b409511e9a462bf7ba03b1cc215
\paragraph{Settled source reading}
\begin{sloppypar}
Calibration uses the standard conditional-\allowbreak{}expectation event identity intended by the source proof.\allowbreak{}
\end{sloppypar}

\paragraph{Lean-expanded paper semantic target}
\begin{ReviewVerbatim}
type:
{X : Type u_1} →
  {Y : Type u_2} →
    [inst : MeasurableSpace (X × Y)] → [DecidableEq Y] → MeasureTheory.Measure (X × Y) → (X → Y → ℝ) → Prop

value:
fun {X : Type u_1} {Y : Type u_2} [MeasurableSpace (X × Y)] [DecidableEq Y] (μ : MeasureTheory.Measure (X × Y))
    (q : X → Y → ℝ) =>
  ∀ (y : Y) (s : Set ℝ),
    MeasurableSet s →
      ∫ (xy : X × Y), if q xy.1 y ∈ s then if xy.2 = y then 1 else 0 else 0 ∂μ =
        ∫ (xy : X × Y), if q xy.1 y ∈ s then q xy.1 y else 0 ∂μ
\end{ReviewVerbatim}


\paragraph{Recorded source-to-declaration screening}
\textbf{Verdict:} \texttt{matches}
\\
{\raggedright
\textbf{Reason:} The approved source reading requires, for every label y and measurable score-\allowbreak{}value event A, equality between the mass of observations with Y = y and q\_\allowbreak{}y in A and the integral of q\_\allowbreak{}y over that same event.\allowbreak{} The expanded calibrated predicate states exactly this event-\allowbreak{}preimage integral identity for every y and measurable set s.\allowbreak{} Its displayed use in the Theorem 1(iii) Spec adds posteriorSimplex and coordinate measurability before assuming calibrated, so the paper use remains within the approved probabilistic score domain.\allowbreak{}
\par}
\reviewmatch{Matches source input}
\noindent\textbf{Reviewer annotation}\par
\reviewerbox
\clearpage
\hypertarget{paper-prerequisite-46e579d6fcd53814}{}
\subsection*{\small\ttfamily\raggedright DSWG24\allowbreak{}Discretization\allowbreak{}Bias.\allowbreak{}Proof\allowbreak{}Bridge.\allowbreak{}is\allowbreak{}Aggregate\allowbreak{}First\allowbreak{}Argmax}
\noindent\textbf{Paper-local declaration:}\par
{\footnotesize\ttfamily\raggedright DSWG24\allowbreak{}Discretization\allowbreak{}Bias.\allowbreak{}Proof\allowbreak{}Bridge.\allowbreak{}is\allowbreak{}Aggregate\allowbreak{}First\allowbreak{}Argmax\par}
\textbf{Source locator:} {\footnotesize\raggedright cited publication, lines 749-\allowbreak{}755\par}
\paragraph{Verbatim paper-source connection}
\begin{ReviewVerbatim}
q
    For the result, we define the notion of non-trivial reference distributions PN     : a reference distribution
                                                       1
pq : {xi } → △NY is in P N if it assigns mass at least N to the most likely class in the  dataset, i.e., at least 1
data point would be discretized to the most likely class arg maxy i q(y, xi ).13
                                                                   P


[Next verbatim source excerpt]

Aggregate posterior reference distribution What is an appropriate reference distribution pref ? Of
course, if we knew the true distribution of labels (e.g., the true fraction of each group in the voter file), we
could measure and optimize fidelity with respect to it. However, this information is not often known. In this
paper, we will thus primarily compare to the aggregate posterior: for a given set {xi }, what would the
decision distribution be if we could make continuous decisions corresponding to the continuous classifier q,
                                                                  1 X
                                            pqagg (y, {xi }) =        q(y, xi ).
                                                                  N i
\end{ReviewVerbatim}



\paragraph{Lean-expanded paper semantic target}
\begin{ReviewVerbatim}
type:
{N K : ℕ} → (Fin N → Fin K → ℝ) → Fin K → Prop

value:
fun {N K : ℕ} (posterior : Fin N → Fin K → ℝ) (target : Fin K) =>
  DSWG24DiscretizationBias.ProofBridge.isFirstArgmax (DSWG24DiscretizationBias.ProofBridge.aggregatePosterior posterior)
    target
\end{ReviewVerbatim}

\paragraph{Direct retained dependencies}
\begin{itemize}\raggedright
\item \hyperlink{governing-declaration-71a4a73d8c8948eb}{DSWG24\allowbreak{}Discretization\allowbreak{}Bias.\allowbreak{}Proof\allowbreak{}Bridge.\allowbreak{}aggregate\allowbreak{}Posterior}
\item \hyperlink{governing-declaration-af44c4b45082ddcb}{DSWG24\allowbreak{}Discretization\allowbreak{}Bias.\allowbreak{}Proof\allowbreak{}Bridge.\allowbreak{}is\allowbreak{}First\allowbreak{}Argmax}
\end{itemize}
\paragraph{Recorded source-to-declaration screening}
\textbf{Verdict:} \texttt{matches}
\\
{\raggedright
\textbf{Reason:} The source's dataset-\allowbreak{}most-\allowbreak{}likely class maximizes the sum of row posteriors, with a consistent tie order;\allowbreak{} the expanded target takes the first maximizer of the aggregate posterior, and division by positive dataset size does not change that argmax in the displayed N>K uses.\allowbreak{}
\par}
\reviewmatch{Matches source input}
\noindent\textbf{Reviewer annotation}\par
\reviewerbox
\clearpage
\hypertarget{paper-prerequisite-2e16c9b0a6f0db1c}{}
\subsection*{\small\ttfamily\raggedright DSWG24\allowbreak{}Discretization\allowbreak{}Bias.\allowbreak{}Proof\allowbreak{}Bridge.\allowbreak{}is\allowbreak{}Argmax\allowbreak{}Rule}
\noindent\textbf{Paper-local declaration:}\par
{\footnotesize\ttfamily\raggedright DSWG24\allowbreak{}Discretization\allowbreak{}Bias.\allowbreak{}Proof\allowbreak{}Bridge.\allowbreak{}is\allowbreak{}Argmax\allowbreak{}Rule\par}
\textbf{Source locator:} {\footnotesize\raggedright cited publication, lines 218\par}
\paragraph{Verbatim paper-source connection}
\begin{ReviewVerbatim}
Argmax rule assigns, for each data point, the most likely2 class: Dargmax ({q(xi }) ≜ {ŷi = arg maxy q(y, xi )}.

[Next verbatim source excerpt]

   2 Suppose ties are broken according to some consistent ordering over the classes.      In the event of ties, we will use
arg maxy q(y, x) to denote the single class that is in the argmax set and is first in the consistent ordering among classes
in the argmax set.
\end{ReviewVerbatim}



\paragraph{Lean-expanded paper semantic target}
\begin{ReviewVerbatim}
type:
{X : Type u_1} → {Y : Type u_2} → (X → Y → ℝ) → (X → Y) → Prop

value:
fun {X : Type u_1} {Y : Type u_2} (q : X → Y → ℝ) (rule : X → Y) =>
  AppliedModelingLib.Learning.Prediction.IsArgmaxRule q rule
\end{ReviewVerbatim}

\paragraph{Direct retained dependencies}
\begin{itemize}\raggedright
\item \hyperlink{governing-declaration-d9f4aa229cbce6df}{Applied\allowbreak{}Modeling\allowbreak{}Lib.\allowbreak{}Learning.\allowbreak{}Prediction.\allowbreak{}Is\allowbreak{}Argmax\allowbreak{}Rule}
\end{itemize}
\paragraph{Recorded source-to-declaration screening}
\textbf{Verdict:} \texttt{matches}
\\
{\raggedright
\textbf{Reason:} The expanded predicate requires the selected label to attain the row-\allowbreak{}wise maximum.\allowbreak{} It relaxes the source's fixed tie ordering, but every source tie-\allowbreak{}broken argmax satisfies it and the displayed theorem1iii use therefore proves a valid strengthening covering every source-\allowbreak{}admissible rule without changing behavior on that domain.\allowbreak{}
\par}
\reviewmatch{Matches source input}
\noindent\textbf{Reviewer annotation}\par
\reviewerbox
\clearpage
\hypertarget{paper-prerequisite-536b545abe995da0}{}
\subsection*{\small\ttfamily\raggedright DSWG24\allowbreak{}Discretization\allowbreak{}Bias.\allowbreak{}Proof\allowbreak{}Bridge.\allowbreak{}is\allowbreak{}Independent\allowbreak{}Randomized\allowbreak{}Rule}
\noindent\textbf{Paper-local declaration:}\par
{\footnotesize\ttfamily\raggedright DSWG24\allowbreak{}Discretization\allowbreak{}Bias.\allowbreak{}Proof\allowbreak{}Bridge.\allowbreak{}is\allowbreak{}Independent\allowbreak{}Randomized\allowbreak{}Rule\par}
\textbf{Source locator:} {\footnotesize\raggedright cited publication, lines 615-\allowbreak{}618\par}
\paragraph{Verbatim paper-source connection}
\begin{ReviewVerbatim}
Decision (discretization) rules are as defined in Section 1.1. Our theoretical results distinguish between
independent (e.g., argmax, threhsolding, Thompson sampling) and joint (e.g., our optimization approach)
rules. Formally, with independent rules, there exists a (potentially randomized) function d such that
Dd ({q(xi )}) = {ŷi = d(q(y, xi ))}.
\end{ReviewVerbatim}



\paragraph{Lean-expanded paper semantic target}
\begin{ReviewVerbatim}
type:
{X : Type u_1} → {N K : ℕ} → ((Fin N → X) → (Fin N → Fin K) → ℝ) → Prop

value:
fun {X : Type u_1} {N K : ℕ} (jointProbability : (Fin N → X) → (Fin N → Fin K) → ℝ) =>
  ∃ (rowProbability : X → Fin K → ℝ),
    (∀ (x : X) (y : Fin K), 0 ≤ rowProbability x y) ∧
      (∀ (x : X), ∑ y : Fin K, rowProbability x y = 1) ∧
        ∀ (xs : Fin N → X) (decision : Fin N → Fin K),
          jointProbability xs decision = ∏ i : Fin N, rowProbability (xs i) (decision i)
\end{ReviewVerbatim}


\paragraph{Recorded source-to-declaration screening}
\textbf{Verdict:} \texttt{matches}
\\
{\raggedright
\textbf{Reason:} The source defines an independent potentially randomized rule by one common row-\allowbreak{}level randomized function d;\allowbreak{} the expanded target represents d as a nonnegative normalized row kernel and factors the conditional joint decision probability into the product of its row probabilities.\allowbreak{}
\par}
\reviewmatch{Matches source input}
\noindent\textbf{Reviewer annotation}\par
\reviewerbox
\clearpage
\hypertarget{paper-prerequisite-45b940eee9bfbd29}{}
\subsection*{\small\ttfamily\raggedright DSWG24\allowbreak{}Discretization\allowbreak{}Bias.\allowbreak{}Proof\allowbreak{}Bridge.\allowbreak{}is\allowbreak{}Independent\allowbreak{}Rule}
\noindent\textbf{Paper-local declaration:}\par
{\footnotesize\ttfamily\raggedright DSWG24\allowbreak{}Discretization\allowbreak{}Bias.\allowbreak{}Proof\allowbreak{}Bridge.\allowbreak{}is\allowbreak{}Independent\allowbreak{}Rule\par}
\textbf{Source locator:} {\footnotesize\raggedright cited publication, lines 225-\allowbreak{}228\par}
\paragraph{Verbatim paper-source connection}
\begin{ReviewVerbatim}
These rules are all independent: the assigned label for a data point i depends on q(xi ) but not on other
data points or their probability vectors {q(xj )}j̸=i .
    In contrast, our proposed approach also includes joint decision optimization rules, which depend on the
entire set of data points, as a solution to discretization bias.

[Next verbatim source excerpt]

data point xi . A (potentially randomized) decision (discretization) rule D : △NY →Y
                                                                                      N
                                                                                         is a rule that, given a
                                   N
set of probability vectors {q(xi )}i=1 associated with N data points, assigns each data point i a label ŷi ∈ Y.

[Next verbatim source excerpt]

    Decision (discretization) rules are as defined in Section 1.1. Our theoretical results distinguish between
independent (e.g., argmax, threhsolding, Thompson sampling) and joint (e.g., our optimization approach)
rules. Formally, with independent rules, there exists a (potentially randomized) function d such that
Dd ({q(xi )}) = {ŷi = d(q(y, xi ))}.
\end{ReviewVerbatim}



\paragraph{Lean-expanded paper semantic target}
\begin{ReviewVerbatim}
type:
{X : Type u_1} → {N K : ℕ} → ((Fin N → X) → Fin N → Fin K) → Prop

value:
fun {X : Type u_1} {N K : ℕ} (rule : (Fin N → X) → Fin N → Fin K) =>
  ∃ (d : X → Fin K), ∀ (xs : Fin N → X) (i : Fin N), rule xs i = d (xs i)
\end{ReviewVerbatim}


\paragraph{Recorded source-to-declaration screening}
\textbf{Verdict:} \texttt{matches}
\\
{\raggedright
\textbf{Reason:} The source says each assigned label depends only on its own row probability vector through one common function d;\allowbreak{} the expanded deterministic target is exactly the existence of d with rule(xs,i)=d(xs i) for every dataset and row.\allowbreak{}
\par}
\reviewmatch{Matches source input}
\noindent\textbf{Reviewer annotation}\par
\reviewerbox
\clearpage
\hypertarget{paper-prerequisite-29d817bb87f95d82}{}
\subsection*{\small\ttfamily\raggedright DSWG24\allowbreak{}Discretization\allowbreak{}Bias.\allowbreak{}Proof\allowbreak{}Bridge.\allowbreak{}is\allowbreak{}Thompson\allowbreak{}Sampling\allowbreak{}Rule}
\noindent\textbf{Paper-local declaration:}\par
{\footnotesize\ttfamily\raggedright DSWG24\allowbreak{}Discretization\allowbreak{}Bias.\allowbreak{}Proof\allowbreak{}Bridge.\allowbreak{}is\allowbreak{}Thompson\allowbreak{}Sampling\allowbreak{}Rule\par}
\textbf{Source locator:} {\footnotesize\raggedright cited publication, lines 222-\allowbreak{}223\par}
\paragraph{Verbatim paper-source connection}
\begin{ReviewVerbatim}
Thompson sampling for each data point i, samples a class y based on the probabilities q(y, xi ), i.e., ŷi = y
    with probability q(y, xi ), and so in expectation will have labels matching the aggregate posterior.
\end{ReviewVerbatim}



\paragraph{Lean-expanded paper semantic target}
\begin{ReviewVerbatim}
type:
{N K : ℕ} → (Fin N → Fin K → ℝ) → (Fin N → Fin K → ℝ) → Prop

value:
fun {N K : ℕ} (q selectionProbability : Fin N → Fin K → ℝ) =>
  ∀ (i : Fin N) (y : Fin K), selectionProbability i y = q i y
\end{ReviewVerbatim}


\paragraph{Recorded source-to-declaration screening}
\textbf{Verdict:} \texttt{matches}
\\
{\raggedright
\textbf{Reason:} The source defines Thompson sampling by selecting label y for row i with probability q(y,x\_\allowbreak{}i);\allowbreak{} the expanded target equates every row-\allowbreak{}label selection probability with the corresponding posterior probability.\allowbreak{}
\par}
\reviewmatch{Matches source input}
\noindent\textbf{Reviewer annotation}\par
\reviewerbox
\clearpage
\hypertarget{paper-prerequisite-f1caeaac865e715b}{}
\subsection*{\small\ttfamily\raggedright DSWG24\allowbreak{}Discretization\allowbreak{}Bias.\allowbreak{}Proof\allowbreak{}Bridge.\allowbreak{}is\allowbreak{}Tie\allowbreak{}Broken\allowbreak{}Argmax\allowbreak{}Rule}
\noindent\textbf{Paper-local declaration:}\par
{\footnotesize\ttfamily\raggedright DSWG24\allowbreak{}Discretization\allowbreak{}Bias.\allowbreak{}Proof\allowbreak{}Bridge.\allowbreak{}is\allowbreak{}Tie\allowbreak{}Broken\allowbreak{}Argmax\allowbreak{}Rule\par}
\textbf{Source locator:} {\footnotesize\raggedright cited publication, lines 215-\allowbreak{}231\par}
\paragraph{Verbatim paper-source connection}
\begin{ReviewVerbatim}
1.1.2    Overview of decision rules
Several rules are used in academia and in practice:

Argmax rule assigns, for each data point, the most likely2 class: Dargmax ({q(xi }) ≜ {ŷi = arg maxy q(y, xi )}.
Threshold at t rule assigns a label only if the argmax class has a probability of at least t, i.e., ŷi =
    arg maxy q(y, xi ) if maxy q(y, xi ) ≥ t, otherwise i is left uncoded.

Thompson sampling for each data point i, samples a class y based on the probabilities q(y, xi ), i.e., ŷi = y
    with probability q(y, xi ), and so in expectation will have labels matching the aggregate posterior.

    These rules are all independent: the assigned label for a data point i depends on q(xi ) but not on other
data points or their probability vectors {q(xj )}j̸=i .
    In contrast, our proposed approach also includes joint decision optimization rules, which depend on the
entire set of data points, as a solution to discretization bias.
   2 Suppose ties are broken according to some consistent ordering over the classes.      In the event of ties, we will use
arg maxy q(y, x) to denote the single class that is in the argmax set and is first in the consistent ordering among classes
in the argmax set.

[Next verbatim source excerpt]

Argmax rule assigns, for each data point, the most likely2 class: Dargmax ({q(xi }) ≜ {ŷi = arg maxy q(y, xi )}.
Threshold at t rule assigns a label only if the argmax class has a probability of at least t, i.e., ŷi =
    arg maxy q(y, xi ) if maxy q(y, xi ) ≥ t, otherwise i is left uncoded.

Thompson sampling for each data point i, samples a class y based on the probabilities q(y, xi ), i.e., ŷi = y
    with probability q(y, xi ), and so in expectation will have labels matching the aggregate posterior.

    These rules are all independent: the assigned label for a data point i depends on q(xi ) but not on other
data points or their probability vectors {q(xj )}j̸=i .
    In contrast, our proposed approach also includes joint decision optimization rules, which depend on the
entire set of data points, as a solution to discretization bias.
   2 Suppose ties are broken according to some consistent ordering over the classes.      In the event of ties, we will use
arg maxy q(y, x) to denote the single class that is in the argmax set and is first in the consistent ordering among classes
in the argmax set.
\end{ReviewVerbatim}



\paragraph{Lean-expanded paper semantic target}
\begin{ReviewVerbatim}
type:
{N K : ℕ} → (Fin N → Fin K → ℝ) → (Fin N → Fin K) → Prop

value:
fun {N K : ℕ} (q : Fin N → Fin K → ℝ) (rule : Fin N → Fin K) =>
  ∀ (i : Fin N), DSWG24DiscretizationBias.ProofBridge.isFirstArgmax (q i) (rule i)
\end{ReviewVerbatim}

\paragraph{Direct retained dependencies}
\begin{itemize}\raggedright
\item \hyperlink{governing-declaration-af44c4b45082ddcb}{DSWG24\allowbreak{}Discretization\allowbreak{}Bias.\allowbreak{}Proof\allowbreak{}Bridge.\allowbreak{}is\allowbreak{}First\allowbreak{}Argmax}
\end{itemize}
\paragraph{Recorded source-to-declaration screening}
\textbf{Verdict:} \texttt{matches}
\\
{\raggedright
\textbf{Reason:} The source chooses a row-\allowbreak{}wise maximum and resolves ties by a fixed consistent class ordering;\allowbreak{} the expanded target uses the natural Fin ordering, requires maximal score, and requires the chosen index to be first among all tied maxima for every row.\allowbreak{}
\par}
\reviewmatch{Matches source input}
\noindent\textbf{Reviewer annotation}\par
\reviewerbox
\clearpage
\hypertarget{paper-prerequisite-ce0ba7b00b93c0fc}{}
\subsection*{\small\ttfamily\raggedright DSWG24\allowbreak{}Discretization\allowbreak{}Bias.\allowbreak{}Proof\allowbreak{}Bridge.\allowbreak{}maximizes\allowbreak{}Equation1}
\noindent\textbf{Paper-local declaration:}\par
{\footnotesize\ttfamily\raggedright DSWG24\allowbreak{}Discretization\allowbreak{}Bias.\allowbreak{}Proof\allowbreak{}Bridge.\allowbreak{}maximizes\allowbreak{}Equation1\par}
\textbf{Source locator:} {\footnotesize\raggedright cited publication, lines 252-\allowbreak{}264\par}
\paragraph{Verbatim paper-source connection}
\begin{ReviewVerbatim}
Integer optimization rules directly optimize given objectives. We consider the rule that corresponds to
     solving the following optimization problem, for a given γ, data points {xi }, and reference pref :
                                                            N
                                                                           !
                      γ                                  1 X
                   D ({q(xi )}, pref ) = arg max{ŷi } γ       q(ŷi , xi ) + (1 − γ)fid(pref , ŷ1:N ). (1)
                                                         N i=1

     The rule balances, parameterized by γ ∈ [0, 1], row-level scores (accuracy as measured by q) and
     distributional fidelity.
Aggregate posterior matching A special case of the joint optimization decision rules defined in Equa-
    tion (1) is as γ → 0, leading to a rule that maximizes accuracy subject to the constraint that distri-
    butional fidelity be as high as possible. We call such rules matching solutions, e.g., aggregate posterior

[Next verbatim source excerpt]


[form-feed in source extraction]
    (2a) To measure bias, we use the distance between the marginal distribution of labels and a reference
distribution pref . Let p̂marg (y) denote the fraction of datapoints with the label y:
                                                                     N
                                                                 1 X
                                          p̂marg (y, ŷ1:N ) =         1[ŷi = y].
                                                                 N i=1

Then, bias is how much a class y is amplified by labels relative to the reference:

                                    bias(y, ŷ1:N , pref ) ≜ p̂marg (y, ŷ1:N ) − pref (y).

    (2b) As a summary statistic, we will further consider distributional fidelity, the negative sum of the
absolute value of bias across classes (equivalently, the negative of the ell1 distance between the label and
reference distributions).
                                                        X
                                 fid(pref , ŷ1:N ) ≜ −   bias(y, ŷ1:N , pref ) .
                                                            y∈Y
\end{ReviewVerbatim}



\paragraph{Lean-expanded paper semantic target}
\begin{ReviewVerbatim}
type:
{N K : ℕ} → ℝ → (Fin N → Fin K → ℝ) → (Fin K → ℝ) → (Fin N → Fin K) → Prop

value:
fun {N K : ℕ} (γ : ℝ) (q : Fin N → Fin K → ℝ) (pref : Fin K → ℝ) (decision : Fin N → Fin K) =>
  ∀ (other : Fin N → Fin K),
    DSWG24DiscretizationBias.ProofBridge.equation1Objective γ q pref other ≤
      DSWG24DiscretizationBias.ProofBridge.equation1Objective γ q pref decision
\end{ReviewVerbatim}

\paragraph{Direct retained dependencies}
\begin{itemize}\raggedright
\item \hyperlink{governing-declaration-78f0292bcb487fbc}{DSWG24\allowbreak{}Discretization\allowbreak{}Bias.\allowbreak{}Proof\allowbreak{}Bridge.\allowbreak{}equation1\allowbreak{}Objective}
\end{itemize}
\paragraph{Recorded source-to-declaration screening}
\textbf{Verdict:} \texttt{matches}
\\
{\raggedright
\textbf{Reason:} Equation (1) maximizes gamma times average selected posterior score plus one-\allowbreak{}minus-\allowbreak{}gamma times fidelity, where fidelity is the negative L1 bias;\allowbreak{} the expanded dependencies implement those formulas and the target compares the chosen decision against every alternative with the correct maximizing inequality.\allowbreak{}
\par}
\reviewmatch{Matches source input}
\noindent\textbf{Reviewer annotation}\par
\reviewerbox
\clearpage
\hypertarget{paper-prerequisite-b11323803f325211}{}
\subsection*{\small\ttfamily\raggedright DSWG24\allowbreak{}Discretization\allowbreak{}Bias.\allowbreak{}Proof\allowbreak{}Bridge.\allowbreak{}posterior\allowbreak{}Simplex}
\noindent\textbf{Paper-local declaration:}\par
{\footnotesize\ttfamily\raggedright DSWG24\allowbreak{}Discretization\allowbreak{}Bias.\allowbreak{}Proof\allowbreak{}Bridge.\allowbreak{}posterior\allowbreak{}Simplex\par}
\textbf{Source locator:} {\footnotesize\raggedright cited publication, lines 597-\allowbreak{}607\par}
\paragraph{Verbatim paper-source connection}
\begin{ReviewVerbatim}
3.1.1    Continuous classifier and discretization rules
Suppose we have a machine learning classifier q (e.g., trained on some other dataset) that makes continuous
probability predictions for each class, given the data. Let q(y, x) denote the probability placed on class y



                                                       10

[form-feed in source extraction]
given data x. For example, the Bayes optimal classifier is
                                               q ∗ (y, x) ≜ Pr(Y = y|X = x).
We will refer to p as the prior and q as our learned posterior. Let the simplex of possible probability vectors
over the support Y be △Y , and let q(x) ∈ △Y be the vector of probabilities (q(y, x))y .
\end{ReviewVerbatim}



\paragraph{Lean-expanded paper semantic target}
\begin{ReviewVerbatim}
type:
{X : Type u_1} → {Y : Type u_2} → [Fintype Y] → (X → Y → ℝ) → Prop

value:
fun {X : Type u_1} {Y : Type u_2} [Fintype Y] (q : X → Y → ℝ) =>
  AppliedModelingLib.Learning.Prediction.IsSimplexValued q
\end{ReviewVerbatim}

\paragraph{Direct retained dependencies}
\begin{itemize}\raggedright
\item \hyperlink{governing-declaration-d1592e772e7c7e3d}{Applied\allowbreak{}Modeling\allowbreak{}Lib.\allowbreak{}Learning.\allowbreak{}Prediction.\allowbreak{}Is\allowbreak{}Simplex\allowbreak{}Valued}
\end{itemize}
\paragraph{Recorded source-to-declaration screening}
\textbf{Verdict:} \texttt{matches}
\\
{\raggedright
\textbf{Reason:} The source places every posterior vector q(x) in the label simplex;\allowbreak{} the expanded IsSimplexValued predicate requires its coordinates to be nonnegative, sum to one, and lie at most one, exactly the probability-\allowbreak{}vector condition.\allowbreak{}
\par}
\reviewmatch{Matches source input}
\noindent\textbf{Reviewer annotation}\par
\reviewerbox
\clearpage
\hypertarget{paper-prerequisite-e9ff5dcccf634407}{}
\subsection*{\small\ttfamily\raggedright DSWG24\allowbreak{}Discretization\allowbreak{}Bias.\allowbreak{}Proof\allowbreak{}Bridge.\allowbreak{}reference\allowbreak{}Simplex\allowbreak{}At}
\noindent\textbf{Paper-local declaration:}\par
{\footnotesize\ttfamily\raggedright DSWG24\allowbreak{}Discretization\allowbreak{}Bias.\allowbreak{}Proof\allowbreak{}Bridge.\allowbreak{}reference\allowbreak{}Simplex\allowbreak{}At\par}
\textbf{Source locator:} {\footnotesize\raggedright cited publication, lines 749-\allowbreak{}755\par}
\paragraph{Verbatim paper-source connection}
\begin{ReviewVerbatim}
q
    For the result, we define the notion of non-trivial reference distributions PN     : a reference distribution
                                                       1
pq : {xi } → △NY is in P N if it assigns mass at least N to the most likely class in the  dataset, i.e., at least 1
data point would be discretized to the most likely class arg maxy i q(y, xi ).13
                                                                   P
\end{ReviewVerbatim}



\paragraph{Lean-expanded paper semantic target}
\begin{ReviewVerbatim}
type:
{Sample : Type u_1} → {K : ℕ} → (Sample → Fin K → ℝ) → Prop

value:
fun {Sample : Type u_1} {K : ℕ} (prefAt : Sample → Fin K → ℝ) =>
  ∀ (sample : Sample), (∀ (y : Fin K), 0 ≤ prefAt sample y) ∧ ∑ y : Fin K, prefAt sample y = 1
\end{ReviewVerbatim}


\paragraph{Recorded source-to-declaration screening}
\textbf{Verdict:} \texttt{matches}
\\
{\raggedright
\textbf{Reason:} The source's dataset-\allowbreak{}indexed reference distribution takes values in the label simplex;\allowbreak{} the expanded target states pointwise nonnegativity and unit total mass, while the displayed sourcePNq use separately supplies the required at-\allowbreak{}least-\allowbreak{}1/\allowbreak{}N mass on the tie-\allowbreak{}broken aggregate-\allowbreak{}posterior maximizer.\allowbreak{}
\par}
\reviewmatch{Matches source input}
\noindent\textbf{Reviewer annotation}\par
\reviewerbox

\clearpage
\hypertarget{source-claim-1966e78967344232}{}
\hypertarget{source-section-f10b5f68e4acf3dc}{}
\section*{Source claims}
\subsection*{Review row 1}
\textbf{Source locator:} {\footnotesize\raggedright cited publication, lines 642-\allowbreak{}666\par}
\paragraph{Verbatim source input}
\begin{ReviewVerbatim}
Theorem 1. Argmax bias depends on predictive uncertainty, i.e., how much information features x provide
about the true class label. Consider calibrated classifier q and the argmax decision rule Dargmax . Let N > K
and consider a reference distribution of either the aggregate posterior or the prior. Then,
  (i). Suppose x provides no information and q(y, x) = Pr(y) for all x, y. Then, amplification bias is
       maximized and all data points are classified as a plurality class:
                                                  (
                                                    1 − Pr(y), if y = arg maxw p(w),
                            BIAS(y, Dargmax ) =
                                                    − Pr(y),      otherwise.
   9 Note that, though related, mean average error of the continuous predictor q(y, x) is not equivalent to expected zero-one

loss of the discrete decisions. The latter, our measure of decision accuracy, also depends on how the continuous predictor is
discretized.                                n                                           o
  10 For all c of non-zero measure: ∀c ∈
                                                 R
                                             c : (x,y) I[q(y, x) = c]f (x, y)d(y, x) > 0 . Further note that calibration implies
Bayesian consistency, i.e., the continuous model cumulatively assigns the proper mass to each class y: EX [q(y, x)] = Pr(y).


                                                              11

[form-feed in source extraction]
                                                                                   (
                                                                                    1,                                                                        if z = y,
  (ii). Suppose the classifier is perfect, i.e., ∀(x, y) ∼ FXY , we have q(z, x) =                                                                                       Then, there
                                                                                    0,                                                                        otherwise.
        is no amplification bias, BIAS(y, Dargmax ) = 0.
\end{ReviewVerbatim}



\paragraph{Expanded PaperInterface specification}
\begin{ReviewVerbatim}
∀ {X : Type u_1} {Y : Type u_2} [inst : Fintype X] [inst_1 : DecidableEq X] [inst_2 : Fintype Y]
  [inst_3 : DecidableEq Y] (μ : PMF (X × Y)) (q : X → Y → ℝ) (rule : X → Y) (y : Y),
  DSWG24DiscretizationBias.ProofBridge.perfectClassifierOnSupport μ q rule →
    DSWG24DiscretizationBias.Finite.paperBias μ rule (DSWG24DiscretizationBias.ProofBridge.prior μ) y = 0 ∧
      DSWG24DiscretizationBias.Finite.paperBias μ rule (DSWG24DiscretizationBias.Finite.paperAggregatePosterior μ q) y =
        0
\end{ReviewVerbatim}

\paragraph{Retained governing declarations}
\begin{itemize}\raggedright
\item \hyperlink{governing-declaration-7bdd5ebec3f67e16}{DSWG24\allowbreak{}Discretization\allowbreak{}Bias.\allowbreak{}Finite.\allowbreak{}paper\allowbreak{}Aggregate\allowbreak{}Posterior}
\item \hyperlink{governing-declaration-5d2200dbd9b4b93d}{DSWG24\allowbreak{}Discretization\allowbreak{}Bias.\allowbreak{}Finite.\allowbreak{}paper\allowbreak{}Bias}
\item \hyperlink{governing-declaration-28e1fcee264c130f}{DSWG24\allowbreak{}Discretization\allowbreak{}Bias.\allowbreak{}Proof\allowbreak{}Bridge.\allowbreak{}perfect\allowbreak{}Classifier\allowbreak{}On\allowbreak{}Support}
\item \hyperlink{governing-declaration-8260e4a6984dc105}{DSWG24\allowbreak{}Discretization\allowbreak{}Bias.\allowbreak{}Proof\allowbreak{}Bridge.\allowbreak{}prior}
\end{itemize}
\paragraph{Lean-elaborated claim roles}\begin{ReviewVerbatim}
1. parameter: Type u_1
2. parameter: Type u_2
3. parameter: Fintype X
4. parameter: DecidableEq X
5. parameter: Fintype Y
6. parameter: DecidableEq Y
7. parameter: PMF (X × Y)
8. parameter: X → Y → ℝ
9. parameter: X → Y
10. parameter: Y
11. assumption: DSWG24DiscretizationBias.ProofBridge.perfectClassifierOnSupport μ q rule
12. conclusion: DSWG24DiscretizationBias.Finite.paperBias μ rule (DSWG24DiscretizationBias.ProofBridge.prior μ) y = 0 ∧
  DSWG24DiscretizationBias.Finite.paperBias μ rule (DSWG24DiscretizationBias.Finite.paperAggregatePosterior μ q) y = 0
\end{ReviewVerbatim}

\paragraph{Lean proof endpoint (built separately)}\begin{ReviewVerbatim}
DSWG24DiscretizationBias.PaperInterface.theorem1ii_perfect_classifier_zero_bias
\end{ReviewVerbatim}

\paragraph{Recorded source-to-Spec screening}
\textbf{Verdict:} \texttt{matches}
\\
{\raggedright
\textbf{Reason:} The support-\allowbreak{}sensitive perfect-\allowbreak{}classifier premise says that the induced decision is the true label and the score vector is one-\allowbreak{}hot on every positive-\allowbreak{}probability joint atom.\allowbreak{} The target then concludes zero amplification bias for every label under both the prior and aggregate-\allowbreak{}posterior references, exactly as in Theorem 1(ii).\allowbreak{}
\par}
\reviewmatch{Matches source input}
\noindent\textbf{Reviewer annotation}\par
\reviewerbox
\clearpage
\hypertarget{source-claim-61c8bcc1f7628e92}{}
\section*{Review row 2}
\textbf{Source locator:} {\footnotesize\raggedright cited publication, lines 756-\allowbreak{}763\par}
\paragraph{Verbatim source input}
\begin{ReviewVerbatim}
Theorem 2. Consider Bayes optimal q, and N > K. For all F :
   (i). For every γ, pref there exists a joint decision-making rule that maximizes Oγ (D, q, pref ) in Equa-
        tion (3).

  (ii). The argmax decision rule Dargmax maximizes O1 (D, q, pref ), i.e., is accuracy maximizing.
 (iii). Dargmax is the only Pareto optimal independent decision rule for non-trivial pref ∈ P. No independent
        decision rule D maximizes objective Oγ for any γ, unless γ = 1 and D agrees with Dargmax with
        probability 1.
\end{ReviewVerbatim}



\paragraph{Expanded PaperInterface specification}
\begin{ReviewVerbatim}
∀ {Z : Type u_1} {X : Type u_2} [inst : Fintype Z] [inst_1 : DecidableEq Z] {N K : ℕ} (ν : PMF Z) (feature : Z → X)
  (posterior : X → Fin K → ℝ) (rule : Z → Fin K) (argmaxRule : X → Fin K) (aggregateArgmax : (Fin N → X) → Fin K)
  (prefAt : (Fin N → X) → Fin K → ℝ),
  2 ≤ K →
    K < N →
      (∀ (x : X), DSWG24DiscretizationBias.ProofBridge.isFirstArgmax (posterior x) (argmaxRule x)) →
        (0 < AppliedModelingLib.pmfProb ν fun (z : Z) => rule z ≠ argmaxRule (feature z)) →
          DSWG24DiscretizationBias.ProofBridge.sourcePNq
              (fun (sample : Fin N → X) => DSWG24DiscretizationBias.Theorem2iii.sampledPosterior posterior sample)
              aggregateArgmax prefAt →
            0 < N →
              ¬DSWG24DiscretizationBias.Pareto.ParetoOptimal
                  (DSWG24DiscretizationBias.Theorem2iii.expectedDatasetAccuracy
                    (DSWG24DiscretizationBias.Theorem2iii.iidSamplePMF ν N) fun (sample : Fin N → Z) =>
                    DSWG24DiscretizationBias.Theorem2iii.sampledPosterior
                      (fun (z : Z) (y : Fin K) => posterior (feature z) y) sample)
                  (DSWG24DiscretizationBias.Theorem2iii.expectedDatasetFidelityAt
                    (DSWG24DiscretizationBias.Theorem2iii.iidSamplePMF ν N) fun (sample : Fin N → Z) =>
                    prefAt fun (i : Fin N) => feature (sample i))
                  fun (sample : Fin N → Z) => DSWG24DiscretizationBias.Theorem2iii.sampledDecision rule sample
\end{ReviewVerbatim}

\paragraph{Retained governing declarations}
\begin{itemize}\raggedright
\item \hyperlink{governing-declaration-9c38d53b9d664221}{Applied\allowbreak{}Modeling\allowbreak{}Lib.\allowbreak{}pmf\allowbreak{}Prob}
\item \hyperlink{governing-declaration-90962900f461a9f9}{DSWG24\allowbreak{}Discretization\allowbreak{}Bias.\allowbreak{}Pareto.\allowbreak{}Pareto\allowbreak{}Optimal}
\item \hyperlink{governing-declaration-af44c4b45082ddcb}{DSWG24\allowbreak{}Discretization\allowbreak{}Bias.\allowbreak{}Proof\allowbreak{}Bridge.\allowbreak{}is\allowbreak{}First\allowbreak{}Argmax}
\item \hyperlink{governing-declaration-9631b3d6800a2bd3}{DSWG24\allowbreak{}Discretization\allowbreak{}Bias.\allowbreak{}Proof\allowbreak{}Bridge.\allowbreak{}source\allowbreak{}P\allowbreak{}Nq}
\item \hyperlink{governing-declaration-daac633422a44729}{DSWG24\allowbreak{}Discretization\allowbreak{}Bias.\allowbreak{}Theorem2iii.\allowbreak{}expected\allowbreak{}Dataset\allowbreak{}Accuracy}
\item \hyperlink{governing-declaration-8158479406f1d885}{DSWG24\allowbreak{}Discretization\allowbreak{}Bias.\allowbreak{}Theorem2iii.\allowbreak{}expected\allowbreak{}Dataset\allowbreak{}Fidelity\allowbreak{}At}
\item \hyperlink{governing-declaration-b906fbb2f231ab40}{DSWG24\allowbreak{}Discretization\allowbreak{}Bias.\allowbreak{}Theorem2iii.\allowbreak{}iid\allowbreak{}Sample\allowbreak{}PMF}
\item \hyperlink{governing-declaration-05d42de5dba2851d}{DSWG24\allowbreak{}Discretization\allowbreak{}Bias.\allowbreak{}Theorem2iii.\allowbreak{}sampled\allowbreak{}Decision}
\item \hyperlink{governing-declaration-457b4184b0081169}{DSWG24\allowbreak{}Discretization\allowbreak{}Bias.\allowbreak{}Theorem2iii.\allowbreak{}sampled\allowbreak{}Posterior}
\end{itemize}
\paragraph{Lean-elaborated claim roles}\begin{ReviewVerbatim}
1. parameter: Type u_1
2. parameter: Type u_2
3. parameter: Fintype Z
4. parameter: DecidableEq Z
5. parameter: ℕ
6. parameter: ℕ
7. parameter: PMF Z
8. parameter: Z → X
9. parameter: X → Fin K → ℝ
10. parameter: Z → Fin K
11. parameter: X → Fin K
12. parameter: (Fin N → X) → Fin K
13. parameter: (Fin N → X) → Fin K → ℝ
14. assumption: 2 ≤ K
15. assumption: K < N
16. assumption: ∀ (x : X), DSWG24DiscretizationBias.ProofBridge.isFirstArgmax (posterior x) (argmaxRule x)
17. assumption: 0 < AppliedModelingLib.pmfProb ν fun (z : Z) => rule z ≠ argmaxRule (feature z)
18. assumption: DSWG24DiscretizationBias.ProofBridge.sourcePNq
  (fun (sample : Fin N → X) => DSWG24DiscretizationBias.Theorem2iii.sampledPosterior posterior sample) aggregateArgmax
  prefAt
19. assumption: 0 < N
20. conclusion: ¬DSWG24DiscretizationBias.Pareto.ParetoOptimal
    (DSWG24DiscretizationBias.Theorem2iii.expectedDatasetAccuracy
      (DSWG24DiscretizationBias.Theorem2iii.iidSamplePMF ν N) fun (sample : Fin N → Z) =>
      DSWG24DiscretizationBias.Theorem2iii.sampledPosterior (fun (z : Z) (y : Fin K) => posterior (feature z) y) sample)
    (DSWG24DiscretizationBias.Theorem2iii.expectedDatasetFidelityAt
      (DSWG24DiscretizationBias.Theorem2iii.iidSamplePMF ν N) fun (sample : Fin N → Z) =>
      prefAt fun (i : Fin N) => feature (sample i))
    fun (sample : Fin N → Z) => DSWG24DiscretizationBias.Theorem2iii.sampledDecision rule sample
\end{ReviewVerbatim}

\paragraph{Lean proof endpoint (built separately)}\begin{ReviewVerbatim}
DSWG24DiscretizationBias.PaperInterface.theorem2iii_randomized_non_argmax_not_pareto
\end{ReviewVerbatim}

\paragraph{Recorded source-to-Spec screening}
\textbf{Verdict:} \texttt{matches}
\\
{\raggedright
\textbf{Reason:} The finite augmented state Z represents a feature draw together with any realized randomization;\allowbreak{} positive disagreement is therefore the source probability-\allowbreak{}one agreement condition in contrapositive form.\allowbreak{} The posterior-\allowbreak{}score expected-\allowbreak{}accuracy expression is the Bayes-\allowbreak{}reduced source metric, and the target retains K>=2, N>K, the nontrivial reference family, and Pareto nonoptimality.\allowbreak{}
\par}
\reviewmatch{Matches source input}
\noindent\textbf{Reviewer annotation}\par
\reviewerbox
\clearpage
\hypertarget{source-claim-1356bc033e8e977b}{}
\section*{Review row 3}
\textbf{Source locator:} {\footnotesize\raggedright cited publication, lines 756-\allowbreak{}763\par}
\paragraph{Verbatim source input}
\begin{ReviewVerbatim}
Theorem 2. Consider Bayes optimal q, and N > K. For all F :
   (i). For every γ, pref there exists a joint decision-making rule that maximizes Oγ (D, q, pref ) in Equa-
        tion (3).

  (ii). The argmax decision rule Dargmax maximizes O1 (D, q, pref ), i.e., is accuracy maximizing.
 (iii). Dargmax is the only Pareto optimal independent decision rule for non-trivial pref ∈ P. No independent
        decision rule D maximizes objective Oγ for any γ, unless γ = 1 and D agrees with Dargmax with
        probability 1.
\end{ReviewVerbatim}



\paragraph{Expanded PaperInterface specification}
\begin{ReviewVerbatim}
∀ {Z : Type u_1} {X : Type u_2} [inst : Fintype Z] [inst_1 : DecidableEq Z] {N K : ℕ} (ν : PMF Z) (feature : Z → X)
  (posterior : X → Fin K → ℝ) (rule : Z → Fin K) (argmaxRule : X → Fin K) (aggregateArgmax : (Fin N → X) → Fin K)
  (prefAt : (Fin N → X) → Fin K → ℝ) (γ : ℝ),
  2 ≤ K →
    K < N →
      (∀ (x : X), DSWG24DiscretizationBias.ProofBridge.isFirstArgmax (posterior x) (argmaxRule x)) →
        (0 < AppliedModelingLib.pmfProb ν fun (z : Z) => rule z ≠ argmaxRule (feature z)) →
          DSWG24DiscretizationBias.ProofBridge.sourcePNq
              (fun (sample : Fin N → X) => DSWG24DiscretizationBias.Theorem2iii.sampledPosterior posterior sample)
              aggregateArgmax prefAt →
            0 ≤ γ →
              γ < 1 →
                0 < N →
                  ¬∀ (other : (Fin N → Z) → Fin N → Fin K),
                      DSWG24DiscretizationBias.Pareto.weightedObjective γ
                          (DSWG24DiscretizationBias.Theorem2iii.expectedDatasetAccuracy
                            (DSWG24DiscretizationBias.Theorem2iii.iidSamplePMF ν N) fun (sample : Fin N → Z) =>
                            DSWG24DiscretizationBias.Theorem2iii.sampledPosterior
                              (fun (z : Z) (y : Fin K) => posterior (feature z) y) sample)
                          (DSWG24DiscretizationBias.Theorem2iii.expectedDatasetFidelityAt
                            (DSWG24DiscretizationBias.Theorem2iii.iidSamplePMF ν N) fun (sample : Fin N → Z) =>
                            prefAt fun (i : Fin N) => feature (sample i))
                          other ≤
                        DSWG24DiscretizationBias.Pareto.weightedObjective γ
                          (DSWG24DiscretizationBias.Theorem2iii.expectedDatasetAccuracy
                            (DSWG24DiscretizationBias.Theorem2iii.iidSamplePMF ν N) fun (sample : Fin N → Z) =>
                            DSWG24DiscretizationBias.Theorem2iii.sampledPosterior
                              (fun (z : Z) (y : Fin K) => posterior (feature z) y) sample)
                          (DSWG24DiscretizationBias.Theorem2iii.expectedDatasetFidelityAt
                            (DSWG24DiscretizationBias.Theorem2iii.iidSamplePMF ν N) fun (sample : Fin N → Z) =>
                            prefAt fun (i : Fin N) => feature (sample i))
                          fun (sample : Fin N → Z) => DSWG24DiscretizationBias.Theorem2iii.sampledDecision rule sample
\end{ReviewVerbatim}

\paragraph{Retained governing declarations}
\begin{itemize}\raggedright
\item \hyperlink{governing-declaration-9c38d53b9d664221}{Applied\allowbreak{}Modeling\allowbreak{}Lib.\allowbreak{}pmf\allowbreak{}Prob}
\item \hyperlink{governing-declaration-d0facd526e2aa2ff}{DSWG24\allowbreak{}Discretization\allowbreak{}Bias.\allowbreak{}Pareto.\allowbreak{}weighted\allowbreak{}Objective}
\item \hyperlink{governing-declaration-af44c4b45082ddcb}{DSWG24\allowbreak{}Discretization\allowbreak{}Bias.\allowbreak{}Proof\allowbreak{}Bridge.\allowbreak{}is\allowbreak{}First\allowbreak{}Argmax}
\item \hyperlink{governing-declaration-9631b3d6800a2bd3}{DSWG24\allowbreak{}Discretization\allowbreak{}Bias.\allowbreak{}Proof\allowbreak{}Bridge.\allowbreak{}source\allowbreak{}P\allowbreak{}Nq}
\item \hyperlink{governing-declaration-daac633422a44729}{DSWG24\allowbreak{}Discretization\allowbreak{}Bias.\allowbreak{}Theorem2iii.\allowbreak{}expected\allowbreak{}Dataset\allowbreak{}Accuracy}
\item \hyperlink{governing-declaration-8158479406f1d885}{DSWG24\allowbreak{}Discretization\allowbreak{}Bias.\allowbreak{}Theorem2iii.\allowbreak{}expected\allowbreak{}Dataset\allowbreak{}Fidelity\allowbreak{}At}
\item \hyperlink{governing-declaration-b906fbb2f231ab40}{DSWG24\allowbreak{}Discretization\allowbreak{}Bias.\allowbreak{}Theorem2iii.\allowbreak{}iid\allowbreak{}Sample\allowbreak{}PMF}
\item \hyperlink{governing-declaration-05d42de5dba2851d}{DSWG24\allowbreak{}Discretization\allowbreak{}Bias.\allowbreak{}Theorem2iii.\allowbreak{}sampled\allowbreak{}Decision}
\item \hyperlink{governing-declaration-457b4184b0081169}{DSWG24\allowbreak{}Discretization\allowbreak{}Bias.\allowbreak{}Theorem2iii.\allowbreak{}sampled\allowbreak{}Posterior}
\end{itemize}
\paragraph{Lean-elaborated claim roles}\begin{ReviewVerbatim}
1. parameter: Type u_1
2. parameter: Type u_2
3. parameter: Fintype Z
4. parameter: DecidableEq Z
5. parameter: ℕ
6. parameter: ℕ
7. parameter: PMF Z
8. parameter: Z → X
9. parameter: X → Fin K → ℝ
10. parameter: Z → Fin K
11. parameter: X → Fin K
12. parameter: (Fin N → X) → Fin K
13. parameter: (Fin N → X) → Fin K → ℝ
14. parameter: ℝ
15. assumption: 2 ≤ K
16. assumption: K < N
17. assumption: ∀ (x : X), DSWG24DiscretizationBias.ProofBridge.isFirstArgmax (posterior x) (argmaxRule x)
18. assumption: 0 < AppliedModelingLib.pmfProb ν fun (z : Z) => rule z ≠ argmaxRule (feature z)
19. assumption: DSWG24DiscretizationBias.ProofBridge.sourcePNq
  (fun (sample : Fin N → X) => DSWG24DiscretizationBias.Theorem2iii.sampledPosterior posterior sample) aggregateArgmax
  prefAt
20. assumption: 0 ≤ γ
21. assumption: γ < 1
22. assumption: 0 < N
23. conclusion: ¬∀ (other : (Fin N → Z) → Fin N → Fin K),
    DSWG24DiscretizationBias.Pareto.weightedObjective γ
        (DSWG24DiscretizationBias.Theorem2iii.expectedDatasetAccuracy
          (DSWG24DiscretizationBias.Theorem2iii.iidSamplePMF ν N) fun (sample : Fin N → Z) =>
          DSWG24DiscretizationBias.Theorem2iii.sampledPosterior (fun (z : Z) (y : Fin K) => posterior (feature z) y)
            sample)
        (DSWG24DiscretizationBias.Theorem2iii.expectedDatasetFidelityAt
          (DSWG24DiscretizationBias.Theorem2iii.iidSamplePMF ν N) fun (sample : Fin N → Z) =>
          prefAt fun (i : Fin N) => feature (sample i))
        other ≤
      DSWG24DiscretizationBias.Pareto.weightedObjective γ
        (DSWG24DiscretizationBias.Theorem2iii.expectedDatasetAccuracy
          (DSWG24DiscretizationBias.Theorem2iii.iidSamplePMF ν N) fun (sample : Fin N → Z) =>
          DSWG24DiscretizationBias.Theorem2iii.sampledPosterior (fun (z : Z) (y : Fin K) => posterior (feature z) y)
            sample)
        (DSWG24DiscretizationBias.Theorem2iii.expectedDatasetFidelityAt
          (DSWG24DiscretizationBias.Theorem2iii.iidSamplePMF ν N) fun (sample : Fin N → Z) =>
          prefAt fun (i : Fin N) => feature (sample i))
        fun (sample : Fin N → Z) => DSWG24DiscretizationBias.Theorem2iii.sampledDecision rule sample
\end{ReviewVerbatim}

\paragraph{Lean proof endpoint (built separately)}\begin{ReviewVerbatim}
DSWG24DiscretizationBias.PaperInterface.theorem2iii_randomized_weighted_objective_maximizer_agrees_argmax
\end{ReviewVerbatim}

\paragraph{Recorded source-to-Spec screening}
\textbf{Verdict:} \texttt{matches}
\\
{\raggedright
\textbf{Reason:} This is the paired Bayes-\allowbreak{}reduced finite randomized endpoint.\allowbreak{} It retains the source regime, positive disagreement, and nontrivial reference condition, and says that the rule cannot maximize the weighted objective for 0<=gamma<1.\allowbreak{} That interval is the accepted clarification of the source's unless gamma=1 statement, not an economic caveat.\allowbreak{}
\par}
\reviewmatch{Matches source input}
\noindent\textbf{Reviewer annotation}\par
\reviewerbox
\clearpage
\hypertarget{source-claim-4662dc8e8bb960d7}{}
\section*{Review row 4}
\textbf{Source locator:} {\footnotesize\raggedright cited publication, lines 744-\allowbreak{}755\par}
\paragraph{Verbatim source input}
\begin{ReviewVerbatim}
3.3     Results: Individual discretization versus joint optimization
Above, we study how the continuous classifier’s properties induce discretization bias, when using the argmax
rule. We now study discrete decision-making: given an imperfect but Bayes optimal classifier, how can
we make unbiased decisions? We show that joint decision-making, as in the program in Equation (1), is
necessary for Pareto optimality.
                                                                                     q
    For the result, we define the notion of non-trivial reference distributions PN     : a reference distribution
                                                       1
pq : {xi } → △NY is in P N if it assigns mass at least N to the most likely class in the  dataset, i.e., at least 1
data point would be discretized to the most likely class arg maxy i q(y, xi ).13
                                                                   P


[Next verbatim source excerpt]

                                                                                     q
    For the result, we define the notion of non-trivial reference distributions PN     : a reference distribution
                                                       1
pq : {xi } → △NY is in P N if it assigns mass at least N to the most likely class in the  dataset, i.e., at least 1
data point would be discretized to the most likely class arg maxy i q(y, xi ).13
                                                                   P
\end{ReviewVerbatim}



\paragraph{Expanded PaperInterface specification}
\begin{ReviewVerbatim}
∀ {Sample : Type u_1} {N K : ℕ} (posterior : Sample → Fin N → Fin K → ℝ) (aggregateArgmax : Sample → Fin K)
  (prefAt : Sample → Fin K → ℝ),
  DSWG24DiscretizationBias.ProofBridge.sourcePNq posterior aggregateArgmax prefAt ↔
    DSWG24DiscretizationBias.ProofBridge.referenceSimplexAt prefAt ∧
      (∀ (sample : Sample),
          DSWG24DiscretizationBias.ProofBridge.isAggregateFirstArgmax (posterior sample) (aggregateArgmax sample)) ∧
        ∀ (sample : Sample), 1 / ↑N ≤ prefAt sample (aggregateArgmax sample)
\end{ReviewVerbatim}

\paragraph{Retained governing declarations}
\begin{itemize}\raggedright
\item \hyperlink{paper-prerequisite-46e579d6fcd53814}{DSWG24\allowbreak{}Discretization\allowbreak{}Bias.\allowbreak{}Proof\allowbreak{}Bridge.\allowbreak{}is\allowbreak{}Aggregate\allowbreak{}First\allowbreak{}Argmax}
\item \hyperlink{paper-prerequisite-e9ff5dcccf634407}{DSWG24\allowbreak{}Discretization\allowbreak{}Bias.\allowbreak{}Proof\allowbreak{}Bridge.\allowbreak{}reference\allowbreak{}Simplex\allowbreak{}At}
\item \hyperlink{governing-declaration-9631b3d6800a2bd3}{DSWG24\allowbreak{}Discretization\allowbreak{}Bias.\allowbreak{}Proof\allowbreak{}Bridge.\allowbreak{}source\allowbreak{}P\allowbreak{}Nq}
\end{itemize}
\paragraph{Lean-elaborated claim roles}\begin{ReviewVerbatim}
1. parameter: Type u_1
2. parameter: ℕ
3. parameter: ℕ
4. parameter: Sample → Fin N → Fin K → ℝ
5. parameter: Sample → Fin K
6. parameter: Sample → Fin K → ℝ
7. conclusion: DSWG24DiscretizationBias.ProofBridge.sourcePNq posterior aggregateArgmax prefAt ↔
  DSWG24DiscretizationBias.ProofBridge.referenceSimplexAt prefAt ∧
    (∀ (sample : Sample),
        DSWG24DiscretizationBias.ProofBridge.isAggregateFirstArgmax (posterior sample) (aggregateArgmax sample)) ∧
      ∀ (sample : Sample), 1 / ↑N ≤ prefAt sample (aggregateArgmax sample)
\end{ReviewVerbatim}

\paragraph{Lean proof endpoint (built separately)}\begin{ReviewVerbatim}
DSWG24DiscretizationBias.PaperInterface.nontrivial_reference_family_definition
\end{ReviewVerbatim}

\paragraph{Recorded source-to-Spec screening}
\textbf{Verdict:} \texttt{matches}
\\
{\raggedright
\textbf{Reason:} The source family is a dataset-\allowbreak{}indexed reference distribution in the label simplex that assigns at least 1/\allowbreak{}N mass to the dataset aggregate-\allowbreak{}posterior maximizer.\allowbreak{} The target states precisely the simplex, fixed-\allowbreak{}order aggregate argmax, and 1/\allowbreak{}N mass conditions for every dataset.\allowbreak{}
\par}
\reviewmatch{Matches source input}
\noindent\textbf{Reviewer annotation}\par
\reviewerbox
\clearpage
\hypertarget{source-claim-861b407dcf52f582}{}
\section*{Review row 5}
\textbf{Source locator:} {\footnotesize\raggedright cited publication, lines 642-\allowbreak{}650\par}
\paragraph{Verbatim source input}
\begin{ReviewVerbatim}
Theorem 1. Argmax bias depends on predictive uncertainty, i.e., how much information features x provide
about the true class label. Consider calibrated classifier q and the argmax decision rule Dargmax . Let N > K
and consider a reference distribution of either the aggregate posterior or the prior. Then,
  (i). Suppose x provides no information and q(y, x) = Pr(y) for all x, y. Then, amplification bias is
       maximized and all data points are classified as a plurality class:
                                                  (
                                                    1 − Pr(y), if y = arg maxw p(w),
                            BIAS(y, Dargmax ) =
                                                    − Pr(y),      otherwise.
\end{ReviewVerbatim}



\paragraph{Expanded PaperInterface specification}
\begin{ReviewVerbatim}
∀ {X : Type u_1} {Y : Type u_2} [inst : Fintype X] [inst_1 : DecidableEq X] [Nonempty X] [inst_3 : Fintype Y]
  [inst_4 : DecidableEq Y] (μ : PMF (X × Y)) (q : X → Y → ℝ) (z y : Y),
  (∀ (x : X) (a : Y), q x a = DSWG24DiscretizationBias.ProofBridge.prior μ a) →
    (∀ (a : Y), DSWG24DiscretizationBias.ProofBridge.prior μ a ≤ DSWG24DiscretizationBias.ProofBridge.prior μ z) →
      (DSWG24DiscretizationBias.Finite.paperBias μ (fun (x : X) => z) (DSWG24DiscretizationBias.ProofBridge.prior μ) y =
          if z = y then 1 - DSWG24DiscretizationBias.ProofBridge.prior μ y
          else -DSWG24DiscretizationBias.ProofBridge.prior μ y) ∧
        DSWG24DiscretizationBias.Finite.paperBias μ (fun (x : X) => z)
            (DSWG24DiscretizationBias.Finite.paperAggregatePosterior μ q) y =
          if z = y then 1 - DSWG24DiscretizationBias.ProofBridge.prior μ y
          else -DSWG24DiscretizationBias.ProofBridge.prior μ y
\end{ReviewVerbatim}

\paragraph{Retained governing declarations}
\begin{itemize}\raggedright
\item \hyperlink{governing-declaration-7bdd5ebec3f67e16}{DSWG24\allowbreak{}Discretization\allowbreak{}Bias.\allowbreak{}Finite.\allowbreak{}paper\allowbreak{}Aggregate\allowbreak{}Posterior}
\item \hyperlink{governing-declaration-5d2200dbd9b4b93d}{DSWG24\allowbreak{}Discretization\allowbreak{}Bias.\allowbreak{}Finite.\allowbreak{}paper\allowbreak{}Bias}
\item \hyperlink{governing-declaration-8260e4a6984dc105}{DSWG24\allowbreak{}Discretization\allowbreak{}Bias.\allowbreak{}Proof\allowbreak{}Bridge.\allowbreak{}prior}
\end{itemize}
\paragraph{Lean-elaborated claim roles}\begin{ReviewVerbatim}
1. parameter: Type u_1
2. parameter: Type u_2
3. parameter: Fintype X
4. parameter: DecidableEq X
5. assumption: Nonempty X
6. parameter: Fintype Y
7. parameter: DecidableEq Y
8. parameter: PMF (X × Y)
9. parameter: X → Y → ℝ
10. parameter: Y
11. parameter: Y
12. assumption: ∀ (x : X) (a : Y), q x a = DSWG24DiscretizationBias.ProofBridge.prior μ a
13. assumption: ∀ (a : Y), DSWG24DiscretizationBias.ProofBridge.prior μ a ≤ DSWG24DiscretizationBias.ProofBridge.prior μ z
14. conclusion: (DSWG24DiscretizationBias.Finite.paperBias μ (fun (x : X) => z) (DSWG24DiscretizationBias.ProofBridge.prior μ) y =
    if z = y then 1 - DSWG24DiscretizationBias.ProofBridge.prior μ y
    else -DSWG24DiscretizationBias.ProofBridge.prior μ y) ∧
  DSWG24DiscretizationBias.Finite.paperBias μ (fun (x : X) => z)
      (DSWG24DiscretizationBias.Finite.paperAggregatePosterior μ q) y =
    if z = y then 1 - DSWG24DiscretizationBias.ProofBridge.prior μ y
    else -DSWG24DiscretizationBias.ProofBridge.prior μ y
\end{ReviewVerbatim}

\paragraph{Lean proof endpoint (built separately)}\begin{ReviewVerbatim}
DSWG24DiscretizationBias.PaperInterface.theorem1i_no_information_bias
\end{ReviewVerbatim}

\paragraph{Recorded source-to-Spec screening}
\textbf{Verdict:} \texttt{matches}
\\
{\raggedright
\textbf{Reason:} The no-\allowbreak{}information and plurality premises expose the source's q(y,x)=Pr(y) and plurality choice.\allowbreak{} The target gives the displayed bias formula for both source-\allowbreak{}permitted reference choices.\allowbreak{} The source's word maximized has no separate formal objective in the pinned statement, so the displayed formula is the complete mathematical conclusion here.\allowbreak{}
\par}
\reviewmatch{Matches source input}
\noindent\textbf{Reviewer annotation}\par
\reviewerbox
\clearpage
\hypertarget{source-claim-d8ad786b7cc58022}{}
\section*{Review row 6}
\textbf{Source locator:} {\footnotesize\raggedright cited publication, lines 667-\allowbreak{}673\par}
\paragraph{Verbatim source input}
\begin{ReviewVerbatim}
(iii). More generally, in the intermediate regime, argmax bias is upper bounded by the predictive error of
        the classifier q:

                                                                                                   BIAS(y, Dargmax ) ≤ MAE(q).                                                         (4)

       The bound is tight: there exist FXY and q such that Equation (4) holds with equality for the plurality
       class.
\end{ReviewVerbatim}

% report-context-presentation-sha256: dc1fa9810ac3a973e2d1a457fa41cd1dc4334b409511e9a462bf7ba03b1cc215
\paragraph{Settled source reading}
\begin{sloppypar}
Calibration uses the standard conditional-\allowbreak{}expectation event identity intended by the source proof.\allowbreak{}
\end{sloppypar}

\paragraph{Expanded PaperInterface specification}
\begin{ReviewVerbatim}
∀ {X : Type u_1} {Y : Type u_2} [inst : MeasurableSpace X] [inst_1 : MeasurableSpace Y] [MeasurableSingletonClass Y]
  [inst_3 : Fintype Y] [inst_4 : DecidableEq Y] (μ : MeasureTheory.Measure (X × Y)) [MeasureTheory.IsFiniteMeasure μ]
  (q : X → Y → ℝ) (argmaxRule : X → Y) (y : Y),
  Measurable argmaxRule →
    DSWG24DiscretizationBias.ProofBridge.isArgmaxRule q argmaxRule →
      DSWG24DiscretizationBias.ProofBridge.posteriorSimplex q →
        (∀ (y : Y), Measurable fun (xy : X × Y) => q xy.1 y) →
          DSWG24DiscretizationBias.ProofBridge.calibrated μ q →
            DSWG24DiscretizationBias.ProofBridge.continuousJointPriorBias μ argmaxRule y ≤
                DSWG24DiscretizationBias.ProofBridge.continuousJointClassifierMAE μ q ∧
              DSWG24DiscretizationBias.ProofBridge.continuousJointAggregateBias μ q argmaxRule y ≤
                DSWG24DiscretizationBias.ProofBridge.continuousJointClassifierMAE μ q
\end{ReviewVerbatim}

\paragraph{Retained governing declarations}
\begin{itemize}\raggedright
\item \hyperlink{paper-prerequisite-c94cbc772446ffec}{DSWG24\allowbreak{}Discretization\allowbreak{}Bias.\allowbreak{}Proof\allowbreak{}Bridge.\allowbreak{}calibrated}
\item \hyperlink{governing-declaration-302610b6fa0b99db}{DSWG24\allowbreak{}Discretization\allowbreak{}Bias.\allowbreak{}Proof\allowbreak{}Bridge.\allowbreak{}continuous\allowbreak{}Joint\allowbreak{}Aggregate\allowbreak{}Bias}
\item \hyperlink{governing-declaration-d07991e4ae99da3b}{DSWG24\allowbreak{}Discretization\allowbreak{}Bias.\allowbreak{}Proof\allowbreak{}Bridge.\allowbreak{}continuous\allowbreak{}Joint\allowbreak{}Classifier\allowbreak{}MAE}
\item \hyperlink{governing-declaration-7df1fb937c21cdf5}{DSWG24\allowbreak{}Discretization\allowbreak{}Bias.\allowbreak{}Proof\allowbreak{}Bridge.\allowbreak{}continuous\allowbreak{}Joint\allowbreak{}Prior\allowbreak{}Bias}
\item \hyperlink{paper-prerequisite-2e16c9b0a6f0db1c}{DSWG24\allowbreak{}Discretization\allowbreak{}Bias.\allowbreak{}Proof\allowbreak{}Bridge.\allowbreak{}is\allowbreak{}Argmax\allowbreak{}Rule}
\item \hyperlink{paper-prerequisite-b11323803f325211}{DSWG24\allowbreak{}Discretization\allowbreak{}Bias.\allowbreak{}Proof\allowbreak{}Bridge.\allowbreak{}posterior\allowbreak{}Simplex}
\end{itemize}
\paragraph{Lean-elaborated claim roles}\begin{ReviewVerbatim}
1. parameter: Type u_1
2. parameter: Type u_2
3. parameter: MeasurableSpace X
4. parameter: MeasurableSpace Y
5. assumption: MeasurableSingletonClass Y
6. parameter: Fintype Y
7. parameter: DecidableEq Y
8. parameter: MeasureTheory.Measure (X × Y)
9. assumption: MeasureTheory.IsFiniteMeasure μ
10. parameter: X → Y → ℝ
11. parameter: X → Y
12. parameter: Y
13. assumption: Measurable argmaxRule
14. assumption: DSWG24DiscretizationBias.ProofBridge.isArgmaxRule q argmaxRule
15. assumption: DSWG24DiscretizationBias.ProofBridge.posteriorSimplex q
16. assumption: ∀ (y : Y), Measurable fun (xy : X × Y) => q xy.1 y
17. assumption: DSWG24DiscretizationBias.ProofBridge.calibrated μ q
18. conclusion: DSWG24DiscretizationBias.ProofBridge.continuousJointPriorBias μ argmaxRule y ≤
    DSWG24DiscretizationBias.ProofBridge.continuousJointClassifierMAE μ q ∧
  DSWG24DiscretizationBias.ProofBridge.continuousJointAggregateBias μ q argmaxRule y ≤
    DSWG24DiscretizationBias.ProofBridge.continuousJointClassifierMAE μ q
\end{ReviewVerbatim}

\paragraph{Lean proof endpoint (built separately)}\begin{ReviewVerbatim}
DSWG24DiscretizationBias.PaperInterface.theorem1iii_argmax_bias_le_mae
\end{ReviewVerbatim}

\paragraph{Recorded source-to-Spec screening}
\textbf{Verdict:} \texttt{matches}
\\
{\raggedright
\textbf{Reason:} The selected source clause bounds each label's argmax bias by classifier MAE for either the prior or aggregate-\allowbreak{}posterior reference.\allowbreak{} The expanded target proves both corresponding inequalities:\allowbreak{} predicted argmax mass minus true-\allowbreak{}label mass and predicted argmax mass minus aggregate posterior are each at most the integral of the simplex score error sum q\_\allowbreak{}y(1-\allowbreak{}q\_\allowbreak{}y).\allowbreak{} Its weak pointwise-\allowbreak{}argmax premise includes the source argmax rule, posteriorSimplex supplies probability-\allowbreak{}vector scores, and the calibrated event identity is exactly the approved usual calibration reading.\allowbreak{} Finiteness and measurability are analytic well-\allowbreak{}formedness restrictions, while allowing any finite measure and any maximizing tie choice strengthens the bound on the source-\allowbreak{}admissible probability-\allowbreak{}distribution case.\allowbreak{}
\par}
\reviewmatch{Matches source input}
\noindent\textbf{Reviewer annotation}\par
\reviewerbox
\clearpage
\hypertarget{source-claim-d6b9878258465151}{}
\section*{Review row 7}
\textbf{Source locator:} {\footnotesize\raggedright cited publication, lines 667-\allowbreak{}673\par}
\paragraph{Verbatim source input}
\begin{ReviewVerbatim}
(iii). More generally, in the intermediate regime, argmax bias is upper bounded by the predictive error of
        the classifier q:

                                                                                                   BIAS(y, Dargmax ) ≤ MAE(q).                                                         (4)

       The bound is tight: there exist FXY and q such that Equation (4) holds with equality for the plurality
       class.
\end{ReviewVerbatim}



\paragraph{Expanded PaperInterface specification}
\begin{ReviewVerbatim}
DSWG24DiscretizationBias.Finite.paperAggregateBias DSWG24DiscretizationBias.Finite.paperTheorem1TightMu
    DSWG24DiscretizationBias.Finite.paperTheorem1TightQ DSWG24DiscretizationBias.Finite.paperTheorem1TightArgmax 0 =
  DSWG24DiscretizationBias.ProofBridge.classifierMAE DSWG24DiscretizationBias.Finite.paperTheorem1TightMu
    DSWG24DiscretizationBias.Finite.paperTheorem1TightQ
\end{ReviewVerbatim}

\paragraph{Retained governing declarations}
\begin{itemize}\raggedright
\item \hyperlink{governing-declaration-f3bd0c0016772e11}{DSWG24\allowbreak{}Discretization\allowbreak{}Bias.\allowbreak{}Finite.\allowbreak{}paper\allowbreak{}Aggregate\allowbreak{}Bias}
\item \hyperlink{governing-declaration-9de2c431ac58b2a3}{DSWG24\allowbreak{}Discretization\allowbreak{}Bias.\allowbreak{}Finite.\allowbreak{}paper\allowbreak{}Theorem1\allowbreak{}Tight\allowbreak{}Argmax}
\item \hyperlink{governing-declaration-d4243b460c2c568f}{DSWG24\allowbreak{}Discretization\allowbreak{}Bias.\allowbreak{}Finite.\allowbreak{}paper\allowbreak{}Theorem1\allowbreak{}Tight\allowbreak{}Mu}
\item \hyperlink{governing-declaration-193c98b93e9de012}{DSWG24\allowbreak{}Discretization\allowbreak{}Bias.\allowbreak{}Finite.\allowbreak{}paper\allowbreak{}Theorem1\allowbreak{}Tight\allowbreak{}Q}
\item \hyperlink{governing-declaration-2e16128ba60f3c18}{DSWG24\allowbreak{}Discretization\allowbreak{}Bias.\allowbreak{}Proof\allowbreak{}Bridge.\allowbreak{}classifier\allowbreak{}MAE}
\end{itemize}
\paragraph{Lean-elaborated claim roles}\begin{ReviewVerbatim}
1. conclusion: DSWG24DiscretizationBias.Finite.paperAggregateBias DSWG24DiscretizationBias.Finite.paperTheorem1TightMu
    DSWG24DiscretizationBias.Finite.paperTheorem1TightQ DSWG24DiscretizationBias.Finite.paperTheorem1TightArgmax 0 =
  DSWG24DiscretizationBias.ProofBridge.classifierMAE DSWG24DiscretizationBias.Finite.paperTheorem1TightMu
    DSWG24DiscretizationBias.Finite.paperTheorem1TightQ
\end{ReviewVerbatim}

\paragraph{Lean proof endpoint (built separately)}\begin{ReviewVerbatim}
DSWG24DiscretizationBias.PaperInterface.theorem1iii_tight_binary_example
\end{ReviewVerbatim}

\paragraph{Recorded source-to-Spec screening}
\textbf{Verdict:} \texttt{matches}
\\
{\raggedright
\textbf{Reason:} The repaired named witness has one feature, two uniformly likely labels, posterior q(0)=q(1)=1/\allowbreak{}2, and the fixed first-\allowbreak{}class tie break.\allowbreak{} It is calibrated at the sole score value, its aggregate-\allowbreak{}posterior bias for class 0 is 1/\allowbreak{}2, and its classifier MAE is 1/\allowbreak{}2, exactly witnessing the source tightness assertion.\allowbreak{}
\par}
\reviewmatch{Matches source input}
\noindent\textbf{Reviewer annotation}\par
\reviewerbox
\clearpage
\hypertarget{source-claim-82ddd45159ec7ca3}{}
\section*{Review row 8}
\textbf{Source locator:} {\footnotesize\raggedright cited publication, lines 756-\allowbreak{}758\par}
\paragraph{Verbatim source input}
\begin{ReviewVerbatim}
Theorem 2. Consider Bayes optimal q, and N > K. For all F :
   (i). For every γ, pref there exists a joint decision-making rule that maximizes Oγ (D, q, pref ) in Equa-
        tion (3).

[Next verbatim source excerpt]

   Now, note that the objective can be decomposed using the tower property, as follows:
 γ
ON (D, q, pref ) = γACCN (D, q) + (1 − γ)FIDN (D, q, pref )
                        [U+0002 control character]                                                                                     [U+0003 control character]
                 = EF{X} EF{Y }|{X} [γacc(yi , D({q(xi )})) + (1 − γ)fid(pref , D({q(xi )})) | {Xi } = {xi }]
                        "            "     N
                                                                                                                      ##
                                        1 X
                 = EF{X} EF{Y }|{X} γ          I [D({q(xi )}) = yi ] + (1 − γ)fid(pref , D({q(xi )})) | {Xi } = {xi }
                                       N i
                        "            "     N
                                                                                                                     ##
                                        1 X
                 = EF{X} EF{Y }|{X} γ          q(D({q(xj )})i , xi ) + (1 − γ)fid(pref , D({q(xi )})) | {Xi } = {xi }
                                       N i

   Where the last line follows from, given Bayes optimal q, we have for all ŷi

                                            q(ŷi , xi ) = EFY |X [Y = ŷi |X = xi ].

Thus, the optimization decision rule maximizes the inner expectation in the final line, for each given dataset
                                                γ
{xi }. Thus, the corresponding rule maximizes ON  (D, q, pref ), as desired.
\end{ReviewVerbatim}



\paragraph{Expanded PaperInterface specification}
\begin{ReviewVerbatim}
∀ {ω : Type u_1} {σ : Type u_2} {N K : ℕ} [inst : NeZero K],
  2 ≤ K →
    K < N →
      ∀ (expect : (ω → ℝ) → ℝ),
        AppliedModelingLib.Decision.FiniteLinearExpectation expect →
          ∀ (observedDataset : ω → σ) (trueLabels : ω → Fin N → Fin K) (γ : ℝ) (posterior : σ → Fin N → Fin K → ℝ)
            (fidelityTerm : σ → (Fin N → Fin K) → ℝ),
            (∀ (i : Fin N) (choose : σ → Fin K),
                (expect fun (x : ω) => if choose (observedDataset x) = trueLabels x i then 1 else 0) =
                  expect fun (x : ω) => posterior (observedDataset x) i (choose (observedDataset x))) →
              ∃ (optRule : σ → Fin N → Fin K),
                ∀ (rule : σ → Fin N → Fin K),
                  DSWG24DiscretizationBias.ProofBridge.expectedObjective expect observedDataset trueLabels γ
                      fidelityTerm rule ≤
                    DSWG24DiscretizationBias.ProofBridge.expectedObjective expect observedDataset trueLabels γ
                      fidelityTerm optRule
\end{ReviewVerbatim}

\paragraph{Retained governing declarations}
\begin{itemize}\raggedright
\item \hyperlink{governing-declaration-443446486047239c}{Applied\allowbreak{}Modeling\allowbreak{}Lib.\allowbreak{}Decision.\allowbreak{}Finite\allowbreak{}Linear\allowbreak{}Expectation}
\item \hyperlink{governing-declaration-8e0892ece8e4e41c}{DSWG24\allowbreak{}Discretization\allowbreak{}Bias.\allowbreak{}Proof\allowbreak{}Bridge.\allowbreak{}expected\allowbreak{}Objective}
\end{itemize}
\paragraph{Lean-elaborated claim roles}\begin{ReviewVerbatim}
1. parameter: Type u_1
2. parameter: Type u_2
3. parameter: ℕ
4. parameter: ℕ
5. assumption: NeZero K
6. assumption: 2 ≤ K
7. assumption: K < N
8. parameter: (ω → ℝ) → ℝ
9. assumption: AppliedModelingLib.Decision.FiniteLinearExpectation expect
10. parameter: ω → σ
11. parameter: ω → Fin N → Fin K
12. parameter: ℝ
13. parameter: σ → Fin N → Fin K → ℝ
14. parameter: σ → (Fin N → Fin K) → ℝ
15. assumption: ∀ (i : Fin N) (choose : σ → Fin K),
  (expect fun (x : ω) => if choose (observedDataset x) = trueLabels x i then 1 else 0) =
    expect fun (x : ω) => posterior (observedDataset x) i (choose (observedDataset x))
16. conclusion: ∃ (optRule : σ → Fin N → Fin K),
  ∀ (rule : σ → Fin N → Fin K),
    DSWG24DiscretizationBias.ProofBridge.expectedObjective expect observedDataset trueLabels γ fidelityTerm rule ≤
      DSWG24DiscretizationBias.ProofBridge.expectedObjective expect observedDataset trueLabels γ fidelityTerm optRule
\end{ReviewVerbatim}

\paragraph{Lean proof endpoint (built separately)}\begin{ReviewVerbatim}
DSWG24DiscretizationBias.PaperInterface.theorem2i_joint_rule_exists
\end{ReviewVerbatim}

\paragraph{Recorded source-to-Spec screening}
\textbf{Verdict:} \texttt{matches}
\\
{\raggedright
\textbf{Reason:} The governing source model at lines 590-\allowbreak{}595 defines |Y| = K >= 2, and Theorem 2 at lines 756-\allowbreak{}758 assumes N > K;\allowbreak{} these are exactly the target premises hK :\allowbreak{} 2 <= K and hNK :\allowbreak{} K < N.\allowbreak{} The Bayes-\allowbreak{}row identity and finite-\allowbreak{}linear expectation encode the tower-\allowbreak{}property conversion used in the displayed proof.\allowbreak{} The target quantifies over a generic gamma and fidelityTerm, so the paper's Equation (3) objective for every stated gamma and pref is a direct specialization, and the existential rule dominates every competing joint rule as claimed.\allowbreak{}
\par}
\reviewmatch{Matches source input}
\noindent\textbf{Reviewer annotation}\par
\reviewerbox
\clearpage
\hypertarget{source-claim-cc8ba41c41752422}{}
\section*{Review row 9}
\textbf{Source locator:} {\footnotesize\raggedright cited publication, lines 760\par}
\paragraph{Verbatim source input}
\begin{ReviewVerbatim}
(ii). The argmax decision rule Dargmax maximizes O1 (D, q, pref ), i.e., is accuracy maximizing.

[Next verbatim source excerpt]

Proof for Part (ii)
    The result follows directly from the classifier being Bayes optimal – for each row, you cannot do better
than picking the most likely class for that row.
    Recall that D({q(xi )}) refers to the set of decisions ŷ1:N with dataset x1:N , where the decisions can be
jointly made across data points. Note that, for the Bayes optimal classifier q ∗ , we have

                                            EFY |X [I [Y = y] | X = x] = q ∗ (y, x)

Thus, we have:
  1
 ON (D, q ∗ , ·) = ACCN (D, q ∗ )                                                                                          (26)
                                        ∗
               = EF [acc(y1:N , D({q (xi )}))]                                                                             (27)
                    "                    #
                      1 X
               = EF          I [ŷi = y]                                                          let ŷi ≜ D({q ∗ (xj )}N
                                                                                                                         j=1 )i
                      N i
                                                                                                                           (28)
                         "             "                                       ##
                                           1 X
               = EF{X} EF{Y }|{X}              I [Yi = ŷi ] | {Xi } = {xi }                                 tower property
                                           N i
                                                                                                                           (29)
                         "                                          #
                             1 X
               = EF{X}           EFY |X [I [Y = ŷi ] | X = xi ]                        independent sampling of datapoints
                             N i
                                                                                                                           (30)
                         "                     #
                             1 X ∗
               = EF{X}           q (ŷi , xi )                                                                             (31)
                             N i
                         "                                  #
                             1 X ∗           ∗
               ≤ EF{X}           q (arg max q (y, xi ), xi )                                                               (32)
                             N i         y

                  1
               ≜ ON (Dargmax , q ∗ , ·).                                                                                   (33)
\end{ReviewVerbatim}



\paragraph{Expanded PaperInterface specification}
\begin{ReviewVerbatim}
∀ {ω : Type u_1} {σ : Type u_2} {N K : ℕ} [NeZero K],
  2 ≤ K →
    K < N →
      ∀ (expect : (ω → ℝ) → ℝ),
        AppliedModelingLib.Decision.FiniteLinearExpectation expect →
          ∀ (observedDataset : ω → σ) (trueLabels : ω → Fin N → Fin K) (posterior : σ → Fin N → Fin K → ℝ)
            {decisionRule argmaxRule : σ → Fin N → Fin K},
            (∀ (xs : σ), AppliedModelingLib.Decision.IsPointwiseMax (posterior xs) (argmaxRule xs)) →
              (∀ (i : Fin N) (choose : σ → Fin K),
                  (expect fun (x : ω) => if choose (observedDataset x) = trueLabels x i then 1 else 0) =
                    expect fun (x : ω) => posterior (observedDataset x) i (choose (observedDataset x))) →
                AppliedModelingLib.Decision.expectedDecisionAccuracy expect observedDataset trueLabels decisionRule ≤
                  AppliedModelingLib.Decision.expectedDecisionAccuracy expect observedDataset trueLabels argmaxRule
\end{ReviewVerbatim}

\paragraph{Retained governing declarations}
\begin{itemize}\raggedright
\item \hyperlink{governing-declaration-443446486047239c}{Applied\allowbreak{}Modeling\allowbreak{}Lib.\allowbreak{}Decision.\allowbreak{}Finite\allowbreak{}Linear\allowbreak{}Expectation}
\item \hyperlink{governing-declaration-06698b1aac1f5804}{Applied\allowbreak{}Modeling\allowbreak{}Lib.\allowbreak{}Decision.\allowbreak{}Is\allowbreak{}Pointwise\allowbreak{}Max}
\item \hyperlink{governing-declaration-61110afe97ed1775}{Applied\allowbreak{}Modeling\allowbreak{}Lib.\allowbreak{}Decision.\allowbreak{}expected\allowbreak{}Decision\allowbreak{}Accuracy}
\end{itemize}
\paragraph{Lean-elaborated claim roles}\begin{ReviewVerbatim}
1. parameter: Type u_1
2. parameter: Type u_2
3. parameter: ℕ
4. parameter: ℕ
5. assumption: NeZero K
6. assumption: 2 ≤ K
7. assumption: K < N
8. parameter: (ω → ℝ) → ℝ
9. assumption: AppliedModelingLib.Decision.FiniteLinearExpectation expect
10. parameter: ω → σ
11. parameter: ω → Fin N → Fin K
12. parameter: σ → Fin N → Fin K → ℝ
13. parameter: σ → Fin N → Fin K
14. parameter: σ → Fin N → Fin K
15. assumption: ∀ (xs : σ), AppliedModelingLib.Decision.IsPointwiseMax (posterior xs) (argmaxRule xs)
16. assumption: ∀ (i : Fin N) (choose : σ → Fin K),
  (expect fun (x : ω) => if choose (observedDataset x) = trueLabels x i then 1 else 0) =
    expect fun (x : ω) => posterior (observedDataset x) i (choose (observedDataset x))
17. conclusion: Real.le✝ (AppliedModelingLib.Decision.expectedDecisionAccuracy expect observedDataset trueLabels decisionRule)
  (AppliedModelingLib.Decision.expectedDecisionAccuracy expect observedDataset trueLabels argmaxRule)
\end{ReviewVerbatim}

\paragraph{Lean proof endpoint (built separately)}\begin{ReviewVerbatim}
DSWG24DiscretizationBias.PaperInterface.theorem2ii_argmax_accuracy_maximizing
\end{ReviewVerbatim}

\paragraph{Recorded source-to-Spec screening}
\textbf{Verdict:} \texttt{matches}
\\
{\raggedright
\textbf{Reason:} The governing source model at lines 590-\allowbreak{}595 defines |Y| = K >= 2, and Theorem 2 at line 756 assumes N > K;\allowbreak{} these are exactly hK :\allowbreak{} 2 <= K and hNK :\allowbreak{} K < N.\allowbreak{} The pointwise-\allowbreak{}max premise and Bayes-\allowbreak{}row identity give the source proof's rowwise posterior comparison, and the conclusion compares the argmax rule against an arbitrary decision rule in expected accuracy, exactly the O\_\allowbreak{}1 claim.\allowbreak{} Source-\allowbreak{}permitted rule randomness is included by instantiating the generic expectation and observation carrier with the rule's random seed;\allowbreak{} the target does not exclude that source case.\allowbreak{}
\par}
\reviewmatch{Matches source input}
\noindent\textbf{Reviewer annotation}\par
\reviewerbox

\end{document}
